\documentclass[11pt,a4paper]{article}
\usepackage[T1]{fontenc}
\usepackage[utf8]{inputenc}
\usepackage{lmodern,amsmath,amssymb,amsthm,mathtools,mathrsfs}
\usepackage{xurl}
\usepackage[margin=28mm]{geometry}
\usepackage{enumitem,microtype,xcolor}
\usepackage[unicode,colorlinks=true,linkcolor=blue!45!black,urlcolor=blue!45!black,citecolor=blue!45!black]{hyperref}
\usepackage{bookmark}
\usepackage{fancyvrb}
\allowdisplaybreaks
\setlength{\emergencystretch}{3em}
\setlist{itemsep=3pt,topsep=5pt}
\numberwithin{equation}{section}
\newtheorem{theorem}{Theorem}[section]
\newtheorem{lemma}[theorem]{Lemma}
\newtheorem{proposition}[theorem]{Proposition}
\newtheorem{corollary}[theorem]{Corollary}
\theoremstyle{definition}
\newtheorem{definition}[theorem]{Definition}
\newtheorem{remark}[theorem]{Remark}
\newcommand{\R}{\mathbb R}
\newcommand{\C}{\mathbb C}
\newcommand{\N}{\mathbb N}
\newcommand{\Z}{\mathbb Z}
\newcommand{\dd}{\,\mathrm d}
\newcommand{\ii}{\mathrm i}
\newcommand{\e}{\mathrm e}
\newcommand{\Rea}{\operatorname{Re}}
\newcommand{\Ima}{\operatorname{Im}}
\newcommand{\supp}{\operatorname{supp}}
\newcommand{\diam}{\operatorname{diam}}
\newcommand{\card}{\operatorname{card}}
\newcommand{\arsinh}{\operatorname{arsinh}}
\newcommand{\src}[2]{\par\smallskip\noindent{\footnotesize\raggedright\textit{Lean correspondence:} \expandafter\nolinkurl\expandafter{\detokenize{#1}}; \expandafter\nolinkurl\expandafter{\detokenize{#2}}.\par}\smallskip}
\newcommand{\lean}[1]{\ifmmode\text{\ttfamily\detokenize{#1}}\else\expandafter\nolinkurl\expandafter{\detokenize{#1}}\fi}
\newcommand{\Mell}{\mathcal M}
\newcommand{\FT}{\mathcal F}
\newcommand{\zabs}{Z}
\newcommand{\densityexp}{p}
\newcommand{\smoothweight}{\psi}
\newcommand{\norm}[1]{\left\lVert#1\right\rVert}
\title{A source-faithful blueprint for a zeta zero-density estimate\\[5pt]
\large The exponent $3(1-\sigma)/(2\sigma)$ and the actual critical-line twelfth moment}
\author{Informal reconstruction of the registered Lean development\\
\small Mathematical development attributed to Mark Lewko and the credited source authors}
\date{Registered source: PALOMAR-2026-10-09-000005, version 1}
\begin{document}
\maketitle
\begin{abstract}
This blueprint reconstructs the mathematical argument implemented in the registered
Lean solution at the immutable source revision specified below. Its two endpoints
are a multiplicity-weighted zero-density estimate for every real $\sigma>3/4$,
and the actual Lebesgue twelfth moment of the Riemann zeta function on $[0,3T]$.
The short-detector argument uses a separately proved Type I large-values estimate;
the long-detector argument uses weighted sampling of the actual critical-line
moment. The moment itself is assembled from an actual fourth moment and a
high-value measure estimate, with the latter produced from a point-to-local-mean
inequality and the signed stationary Atkinson source. All constants, cutoffs,
powers, residue terms, endpoints and quantifier orders are kept visible. Source
names accompany the mathematical statements so that each step can be followed
back into the formal development. Standard analytic foundations supplied by
mathlib are distinguished from analytic results proved in the repository.
\end{abstract}
\tableofcontents
\newpage

\section{Source contract, conventions and the two endpoints}
\label{sec:contract}
The source of this reconstruction is the repository
\begin{center}\small
\url{https://github.com/mrklwk/zeta-density-three-over-two-sigma-lean}
\end{center}
at revision
\begin{center}\small\texttt{804bdb55b1fc40e6a1a6878fd96078644a70e288}.\end{center}
The public endpoints occur in \lean{lean/Solution.lean}. The registry permalink is
\begin{center}\small
\url{https://palomar-registry.org/entry?id=PALOMAR-2026-10-09-000005&version=1}.
\end{center}
The precise machine-readable record is listed in the bibliography. The project's
mathlib revision is \texttt{065356127b1dc0016f66b7283ce0ce2c4055aa55}.
The source snapshot, not the current default branch, governs every correspondence
in this document. Names without an explicit namespace prefix belong to
\lean{MathCollab.Density}, or to its \lean{Stronger} namespace where indicated.
In a source-correspondence list, a bare filename following an explicitly
qualified filename abbreviates that directory, unless the filename uniquely
identifies a different local module (for example \lean{NativePeaks.lean}).
The accompanying source ledger resolves each declaration independently to its
complete path and starting line, so no abbreviated filename needs to be guessed.

\paragraph{Meaning of reconstruction.}
The objects below are ordinary mathematical presentations of the exact Lean
objects. A lemma proved locally is not promoted to an assumption of an endpoint:
when a consumer has an explicit hypothesis, its producer is identified and the
composition is explained. General results about integration, holomorphic
functions, finite sums, real powers and classical arithmetic imported from
mathlib are used with their hypotheses stated. Their full foundational proofs
are not reproduced recursively. Source lemmas dealing with casts, coercions,
finite-support bookkeeping and changes between equivalent integral encodings
are explained with the construction they support, rather than elevated to new
mathematical claims. No assertion about a stronger parameter range, effective
constants, or a sharper logarithmic estimate is intended.

\paragraph{Notation and intervals.}
The function $\zeta$ is mathlib's actual \lean{riemannZeta}. On the pole-free
regions used here it is the usual Riemann zeta function. We use
\[
 s(t)=\tfrac12+\ii t,\qquad \zabs(t)=|\zeta(s(t))|,\qquad d(n)=\#\{a\in\N:a\mid n\}
 \quad(n\ge1),\qquad d(0)=0.
\]
The last convention is important for sums indexed by all natural numbers:
zero-indexed divisor and mollifier terms vanish. For a positive real base,
$x^u=\exp(u\log x)$; for a positive real base regarded as complex,
$x^s=\exp(s\log x)$. All branch-sensitive complex Gaussian powers below are
principal complex powers. An expression $\int_a^b$ is an oriented interval
integral. If $a\le b$, it agrees with the Lebesgue integral over $(a,b]$;
changing a finite number of endpoints does not change it. When a set is displayed,
as in $\int_{[0,3T]}$, the expression is the set-restricted Lebesgue integral.
We explicitly establish integrability wherever the argument requires it.

\paragraph{Eventual assertions and constants.}
``Eventually in $T$'' means that there is a real $T_0$ such that the assertion
holds for every real $T\ge T_0$. When an eventual statement is followed by
``for every $V$, $W$, or $\rho$'', this threshold precedes those variables.
The notation $A\ll_\theta B$ abbreviates $A\le C_\theta B$ for a nonnegative
constant depending only on the displayed fixed parameters. If the Lean producer
gives a positive constant, it can and will be chosen positive here. Constants
may be increased, and a finite collection of thresholds may be replaced by their
maximum; these are exactly the operations used by the eventual-filter proofs.

\begin{theorem}[Registered zero-density endpoint]\label{thm:main}
For every $\sigma,\varepsilon\in\R$ with $\sigma>3/4$ and $\varepsilon>0$, there
exists $C>0$ such that for every real $T\ge2$ there is a finite set
\[
 \mathcal Z_{\sigma,T}=\{\rho\in\C:\zeta(\rho)=0,\ \Rea\rho\ge\sigma,
 |\Ima\rho|\le T\}.
\]
For each $\rho\in\mathcal Z_{\sigma,T}$ the analytic vanishing order
$m(\rho)$ is a positive finite natural number, equal to the extended natural
analytic order, and
\[
 \sum_{\rho\in\mathcal Z_{\sigma,T}}m(\rho)
 \le C T^{3(1-\sigma)/(2\sigma)+\varepsilon}.
\]
The finite set contains every zero satisfying the displayed conditions.
\end{theorem}
\src{Solution.lean}{DensityStronger.density_bound; exact_membership; exact_multiplicity}

\begin{theorem}[Registered actual twelfth moment]\label{thm:moment}
For every $\varepsilon>0$ there exist $K>0$ and $T_0\in\R$ such that for every
real $T\ge T_0$,
\[
 \int_{[0,3T]}\zabs(t)^{12}\dd t\le K T^{2+\varepsilon}.
\]
For every $a,b\in\R$ and every $n\in\N$, the function
$t\mapsto\zabs(t)^n$ is integrable on $[a,b]$.
\end{theorem}
\src{Solution.lean}{DensityStronger.twelfth_moment; actual_moment_integrable}

\paragraph{Restrictions that must not be dropped.}
The density theorem excludes $\sigma=3/4$ and $\varepsilon=0$. Its constant is
chosen after $\sigma,\varepsilon$ and before $T$. It is not asserted uniform as
$\sigma\downarrow3/4$ or $\varepsilon\downarrow0$. Both ordinate signs and both
height boundaries are included. The moment estimate is eventual, whereas the
density estimate holds for all $T\ge2$. The count for $\sigma\ge1$ is empty;
this branch is proved separately and does not require a moment argument.

\section{Actual zeros and their analytic multiplicities}
\label{sec:zeros}
\begin{definition}
A nontrivial zero is a complex number $\rho$ satisfying
\[
 \zeta(\rho)=0,\qquad 0<\Rea\rho<1.
\]
For $\rho\ne1$, let $m(\rho)$ be the natural analytic order of $\zeta$ at
$\rho$. Write
\begin{align*}
 \mathcal Z(\sigma;a,b)&=\{\rho:\rho\text{ is a nontrivial zero},\
  \sigma\le\Rea\rho,\ a\le\Ima\rho\le b\},\\
 N_{\rm slab}(\sigma;a,b)&=\sum_{\rho\in\mathcal Z(\sigma;a,b)}m(\rho),\\
 N(\sigma,T)&=N_{\rm slab}(\sigma;-T,T).
\end{align*}
\end{definition}
\src{ZetaCounting.lean}{IsNontrivialZetaZero; zetaMultiplicity; zetaZeroRegion; zetaZeroFinset; zetaSlabCount; zetaDensityCount}

\begin{lemma}[Finiteness and local factorization]\label{lem:zeroorder}
The set $\mathcal Z(\sigma;a,b)$ is finite for arbitrary real $\sigma,a,b$.
At every nontrivial zero, $0<m(\rho)<\infty$, and there is a function $g$
analytic at $\rho$, with $g(\rho)\ne0$, for which
\[
 \zeta(z)=(z-\rho)^{m(\rho)}g(z)
\]
in a neighborhood of $\rho$. The natural order and the extended natural order
coincide at such a zero.
\end{lemma}
\begin{proof}
Mathlib supplies analyticity of $\zeta$ on $\C\setminus\{1\}$ and finiteness of
its zero set in a compact set. Every zero in the displayed region has
\[
 |\rho|\le|\Rea\rho|+|\Ima\rho|\le1+|a|+|b|.
\]
Thus the region is a subset of the zeros in a compact closed ball, and is finite.

If the analytic order at some $\rho\ne1$ were infinite, the analytic germ would
vanish identically near $\rho$. The identity principle on the connected punctured
plane would then make $\zeta$ identically zero on $\C\setminus\{1\}$, contradicting
$\zeta(2)\ne0$. Hence its extended order is finite. Analyticity and
$\zeta(\rho)=0$ give a nonzero order; conversion from finite extended natural
order to a natural number therefore gives a positive number without loss.
The analytic-order factorization theorem then gives $g$ and the displayed local
identity. A nontrivial zero is not $1$, since its real part is strictly less than
$1$, so all these facts apply.
\end{proof}
\src{ZetaCounting.lean}{zetaZeroRegion_finite; zeta_analyticOrder_ne_top; zetaMultiplicity_pos; zetaMultiplicity_cast; zetaMultiplicity_local_factorization}

\begin{lemma}[Exact membership at the endpoint]
If $\sigma>3/4$, then
\[
 \rho\in\mathcal Z(\sigma;-T,T)
 \iff\zeta(\rho)=0\ \text{and}\ \sigma\le\Rea\rho\ \text{and}\ |\Ima\rho|\le T.
\]
If $\sigma\ge1$, the set is empty for every $T$.
\end{lemma}
\begin{proof}
One implication simply forgets the strip conditions. Conversely, a zero with
$\Rea\rho\ge\sigma>3/4$ has positive real part. Mathlib's nonvanishing theorem
for $\Rea s\ge1$ forces $\Rea\rho<1$. Thus it is a nontrivial zero, and
$-T\le\Ima\rho\le T$ is precisely $|\Ima\rho|\le T$.
If $\sigma\ge1$, membership would give both $\sigma\le\Rea\rho$ and
$\Rea\rho<1$, an impossibility.
\end{proof}
\src{Solution.lean; DensityConclusion.lean}{exact_membership; zetaDensityCount_eq_zero_of_one_le}

\paragraph{Monotonicity and conjugation.}
The slab count increases when either the region is enlarged by lowering the
real-part cutoff, or its lower height is lowered, or its upper height is raised.
This follows by inclusion of the finite zero sets and nonnegativity of
multiplicities. Complex conjugation preserves zeros, real part and analytic
order, and reverses ordinate. Consequently
\[
 N(\sigma,T)\le2N_{\rm slab}(\sigma;0,T).
\]
The inequality, rather than an equality, permits zeros on the real axis and
closed shared boundaries. Similarly
\[
 N_{\rm slab}(\sigma;a,c)\le
 N_{\rm slab}(\sigma;a,b)+N_{\rm slab}(\sigma;b,c)
 \qquad(a\le b\le c).
\]
These finite-set statements are what the final dyadic argument uses.
\src{ZetaCounting.lean; ZetaConjugation.lean}{zetaSlabCount_mono; zetaDensityCount_mono; zetaDensityCount_le_twice_positive; zetaSlabCount_split_le}

\section{Imported analytic foundations and the proved Weyl input}
\label{sec:foundations}
The analytic foundations used throughout are the following, in their exact roles.
\begin{enumerate}
\item Analyticity and the functional equation of actual $\zeta$, complex
conjugation, and nonvanishing on $\Rea s\ge1$.
With $\Gamma_{\R}(s)=\pi^{-s/2}\Gamma(s/2)$, the completed zeta satisfies
$\Lambda(s)=\Lambda(1-s)$ away from the pole conventions. Only pole-free
instances are substituted into quotient identities.
\item The identity theorem and the analytic-order local factorization used in
Section~\ref{sec:zeros}; Jensen's formula for an analytic function on a
neighborhood of a closed disk whose center value is nonzero.
\item Cauchy's integral theorem on rectangles, and simple- and double-pole
residue calculations. The repository proves its rectangle adapters by clearing
the relevant pole numerator, rather than presuming a generic contour shift.
\item Lebesgue dominated convergence, Fubini and Tonelli, comparison of sums
with integrals, and the layer-cake formula for a nonnegative integrable
function. Absolute convergence or a product-space integrable majorant is
established before every infinite interchange.
\item Fourier inversion and Poisson summation for Schwartz functions; compact
smooth functions are made Schwartz by their support and derivative bounds.
The Fourier convention is $\FT f(\xi)=\int_{\R}f(x)e^{-2\pi\ii x\xi}\dd x$.
The Mellin convention is $\Mell f(s)=\int_0^\infty f(x)x^{s-1}\dd x$.
\item The ordinary arithmetic identities for the M\"obius function $\mu$,
Dirichlet convolution and divisor coefficients, together with the elementary
subpower divisor estimate $d(n)\ll_\theta n^\theta$ for every $\theta>0$.
The repository proves the uniform maximum versions from this input.
\item Gaussian integrals, the Gamma functional equation, reflection and
conjugation, Euler's Gamma limit, and elementary differentiation of $\log$,
real powers, $\sqrt{\cdot}$ and $\arsinh$.
\end{enumerate}
These are established mathematical results in the imported library, not new
analytic premises of Theorems~\ref{thm:main} and~\ref{thm:moment}.
The substantial Weyl, fourth-moment, point-mean and Atkinson packet statements
are proved in local source modules and are described below.

\begin{proposition}[Actual critical-line Weyl growth]\label{prop:weyl}
For every $\theta>0$ there is $C_\theta>0$ such that for all real $t$,
\[
 \zabs(t)\le C_\theta(1+|t|)^{1/6+\theta}.
\]
For every $\theta>0$, eventually in $H$, simultaneously for every
$H\le t\le2H$,
\[
 \zabs(t)<H^{1/6+\theta}.
\]
\end{proposition}
\begin{proof}
The local finite Weyl block estimate applies to
\[
 \sum_{n=0}^{B-1}(A+n)^{-1/2-\ii t},\qquad
 t^{1/3}\le A,\quad A^2\le t,\quad B\le A,
\]
and bounds its norm by $30t^{1/6}$. It is obtained from the finite
van der Corput differencing and second-derivative estimates in
\lean{GuthMaynard.Weyl}, with the actual logarithmic phase and the
$1/\sqrt{A+n}$ amplitude retained. If $A\le t^{1/3}$ instead, absolute
estimation gives $B A^{-1/2}\le A^{1/2}\le t^{1/6}$.
A prefix of length $N$ with $N^2\le t$ is partitioned into at most
$\lfloor\log_2N\rfloor+1$ dyadic blocks. The resulting bound is
\[
 \left|\sum_{1\le n\le N}n^{-1/2-\ii t}\right|
 \le30(\log t+1)t^{1/6}.
\]
Fixed multiples of the logarithm are eventually at most $t^\theta$,
uniformly in all such $N$.

The source's sharp approximate functional equation is an identity for actual
$\zeta(s(t))$, with a main sharp sum, its reflected sharp sum multiplied by
$\chi(s(t))$, and a uniformly bounded remainder. Its internally chosen cutoffs
are $x=A+1/2$ and $y=t/(2\pi x)$, with the natural floors $A$ and $M$ satisfying
$A^2,M^2\le t$. The proof of this source interface uses the Abel
sawtooth expansion, Poisson summation, the separate low- and high-frequency
integral estimates, the Gamma factor and the explicit cutoff selection. The
remainder is bounded before $t$ varies; no bound on actual zeta is a hypothesis.
Since $|\chi(s(t))|=1$, the preceding two prefix estimates and the bounded
remainder imply $\zabs(t)<t^{1/6+\theta}$ eventually, after spending a smaller
positive loss on each prefix.

Reflection gives $\zabs(-t)=\zabs(t)$. Continuity on the compact initial
critical-line interval supplies a finite bound there, which is absorbed into
$C_\theta$. This proves the all-real-height statement. For the dyadic statement,
apply the positive-height bound with loss $\theta/2$: for $H\le t\le2H$,
\[
 \zabs(t)<t^{1/6+\theta/2}
 \le2^{1/6+\theta/2}H^{1/6+\theta/2}<H^{1/6+\theta}
\]
once $H^{\theta/2}$ exceeds the fixed factor. All thresholds precede $t$.
\end{proof}
\src{WeylPort/Block.lean; WeylPort/ShortSum.lean; WeylPort/ActualZeta.lean; WeylInput.lean}{norm_criticalLine_cpow_block_le; norm_critical_sum_le_log; eventually_norm_riemannZeta_critical_lt_rpow; exists_riemannZeta_critical_dyadic_weyl; actual_zeta_weyl_input}

\paragraph{Coarse growth used for integrability.}
The Abel continuation proof gives $\zabs(t)\le4(1+|t|)$; a global quadratic
bound for $\zabs(t)^2$ is also proved by combining the high-height Abel bound
with continuity on a compact interval. These coarse bounds suffice for all
polynomially or exponentially weighted integrals. Proposition~\ref{prop:weyl}
implies, with a suitably small loss, that eventually
$\zabs(t)^{12}\le H^3$ throughout $[H,2H]$. This upper truncation is used in
layer cake, not mistaken for the twelfth-moment estimate itself.
\src{Stronger/MomentGrowth.lean; Stronger/PointMean/Native.lean}{zetaMomentCriticalNorm_le_linear; zetaMomentCriticalNorm_twelfth_eventually_le_cube; exists_zetaMomentCriticalNorm_sq_le_quadratic}

\section{Type I large values: objects, reflection and bootstrap}
\label{sec:typeI}
\begin{definition}[Large-value data]
For positive integer $N$, coefficients $a_n\in\C$, and real $t$, put
\[
 D_{N,a}(t)=\sum_{N<n\le2N}a_n e^{\ii t\log(n/N)}.
\]
A finite real set $W$ is one-separated if $|t-u|\ge1$ whenever $t,u\in W$
and $t\ne u$. For fixed $\kappa>0$, admissible data are
$T,V,N,a,W,t_0$ with
\begin{gather*}
 T\ge2,\quad N\ge1,\quad V>0,\quad T^{1/3}\le N\le T,\quad
 N^{3/4+\kappa}\le V,\\
 |a_n|\le1\ (N<n\le2N),\quad W\subset[t_0,t_0+T],\quad
 W\text{ one-separated},\quad |D_{N,a}(t)|\ge V\ (t\in W).
\end{gather*}
No analytic conclusion is a field of this data structure.
For a finite nonempty $U$ set $d_U=1+\max U-\min U$, and put $d_\varnothing=1$.
Finally define
\[
 \Delta(T)=\max_{1\le m\le\lceil32(T+1)\rceil}d(m),\qquad
 b=\frac{N\Delta(T)}{V^2}.
\]
\end{definition}
\src{LargeValueDefinitions.lean}{detectingPolynomial; oneSeparated; localDiameter; divisorMaximum; bootstrapScalar; LargeValueData}

\begin{theorem}[Independent Type I estimate]\label{thm:LV}
For every $\kappa,\varepsilon>0$ there is $C>0$, chosen before every admissible
data set, such that
\[
 |W|\le C T^\varepsilon\left(\frac{N^2}{V^2}
              +\frac{N^3\sqrt T}{V^4}\right).
\]
The starting point $t_0$ is arbitrary.
\end{theorem}
\src{LargeValues.lean}{large_values_data; large_values}

\subsection{The exact comparison objects}
Choose once and for all a nonnegative $C^\infty$ cutoff $w$ equal to $1$ on
$[1,4]$ and supported in $[1/2,5]$. The source constructs it as a smooth bump
centered at $5/2$, inner radius $3/2$, outer radius $2$; its actual outer
support $[1/2,9/2]$ is contained in the stated larger interval. A
\lean{ReflectionCutoff} requires smoothness and support; a
\lean{ComparisonCutoff} additionally requires nonnegativity and the plateau.
Define
\begin{align}
 I_w(v)&=\int_0^\infty w(x)x^{\ii v}\dd x,\\
 h_w(L,v)&=\sum_{n\ge0}w(n/L)(n/L)^{\ii v}-LI_w(v),\\
 P_M(v)&=\sum_{M\le q<2M}(q/M)^{\ii v},\\
 E_M(U)&=\sum_{t,u\in U}|P_M(t-u)|^2,\\
 Q_L(U)&=\sum_{t,u\in U}|h_w(L,t-u)|^2.
\end{align}
The apparently infinite discrete sum defining $h_w$ has finite support, and its
$n=0$ term is zero. Notice the distinct half-open conventions: the auxiliary
polynomial uses $[M,2M)$; the detecting polynomial uses $(N,2N]$.
\src{ReflectionDefinitions.lean; ComparisonCutoff.lean}{ReflectionCutoff; ComparisonCutoff; comparisonCutoff; cutoffIntegral; hKernel; dyadicPolynomial; completePairEnergy; kernelEnergy}

\begin{lemma}[Comparison with the kernel energy]\label{lem:comparison}
There is $C_w\ge0$ such that for all admissible data, every one-separated
subset $U\subset W$, and every positive integer $M$ with $NM\le8(T+1)$,
\[
 E_M(U)\le C_w bNM^2|U|+b^2Q_{NM}(U).
\]
Also $E_M(U)\ge M|U|^2$.
\end{lemma}
\begin{proof}
Expanding the complete pair energy in the $q,r$ variables gives its Gram form
\[
 E_M(U)=\sum_{M\le q,r<2M}
 \left|\sum_{t\in U}(q/r)^{\ii t}\right|^2.
\]
The diagonal $q=r$ contains exactly $M$ copies of $|U|^2$, proving the lower
bound. For the upper bound, insert
$D_{N,a}(t)/D_{N,a}(t)$ inside each $t$ sum. These denominators are nonzero by
$|D_{N,a}(t)|\ge V>0$. Finite Cauchy--Schwarz and $|a_n|\le1$ give the
factor $N$ and sums over products $m=nq$. A fixed product has at most $d(m)$
representations, bounded by $\Delta(T)$ because the scale restriction places it
within the defining maximum. The cutoff equals one on the original normalized
product support, and is nonnegative, so extending to all cutoff-supported
integers does not decrease the sum.

The resulting extension, when its squares are expanded in $t,u$, contains
\[
 NM I_w(t-u)+h_w(NM,t-u)
\]
multiplied by $\overline{P_M(t-u)}$ and divided by
$D_{N,a}(t)\overline{D_{N,a}(u)}$. The continuous piece has
$|I_w(v)|\le C(1+|v|)^{-2}$: set $x=e^{-y}$ and apply Fourier decay to the
compact smooth logarithmic lift. One-separatedness gives a uniform row sum
\[
 \sum_{u\in U}(1+|t-u|)^{-2}\le S
\]
for a fixed finite constant $S=\lean{separationMass}$. Thus the continuous
piece is at most a fixed multiple of $bNM^2|U|$.
For the remaining piece use, pair by pair,
\[
 b|h_w(NM,t-u)|\,|P_M(t-u)|
 \le\tfrac12b^2|h_w(NM,t-u)|^2+\tfrac12|P_M(t-u)|^2.
\]
After summing, the last term is $E_M(U)/2$ and is moved to the left. This
produces the displayed comparison, with its constant chosen from $w$ alone.
\end{proof}
\src{CompletePairs.lean; Comparison.lean}{completePairEnergy_gram; completePairEnergy_diagonal_lower; comparison_extension_bound; comparison_continuous_bound; comparison_remainder_bound; uniform_comparison}

\subsection{Poisson--Mellin reflection, with all frequencies retained}
For $H,L>0$ the retained integer band is exactly
\[
 \mathcal B(L,H)=\left\{m\in\Z:m>0,\quad
 \frac{H}{20\pi L}\le m\le\frac{4H}{\pi L}\right\}.
\]
The mode integral is
\[
 J_w(v,k)=\int_0^\infty w(x)e^{\ii(v\log x-2\pi kx)}\dd x.
\]
Poisson summation gives the absolutely convergent identity
\[
 h_w(L,v)=L\sum_{m\ne0}J_w(v,mL).
\]
For $H\le v\le2H$, the off-band sum has norm at most $C_A L^{1-A}$,
for every integer $A>1$. The same estimate holds for $h_w(L,0)$.
To prove it, the source repeatedly integrates by parts with quotient
$x/(v-2\pi kx)$; off-band the denominator has the explicit lower bound
$(H+|k|)/300$ on $[1/3,6]$. Rescaling puts parameters in a compact set,
so all iterated amplitudes have uniform bounds. Summing
$|m|^{-A}$ justifies the discarded-mode estimate. Endpoint amplitudes vanish
because the support lies strictly within this interval.
\src{UniformIBP.lean; CutoffTransforms.lean; FixedBandError.lean}{ibpIntegral_recurrence; uniform_offBand_decay; hKernel_poisson; uniform_fixedBand_and_diagonal_error}

\begin{lemma}[All-frequency oscillatory estimate]\label{lem:allfreq}
For every real $\tau$ and $0<a\le b$,
\[
 \left|\int_a^b z^{-1}e^{\ii(\tau\log z-2\pi z)}\dd z\right|
 \le\frac6{\sqrt{2\pi a}}.
\]
\end{lemma}
\begin{proof}
Put $s=\sqrt{2\pi a}$ and
$L_0=(\tau-s)/(2\pi)$, $R_0=(\tau+s)/(2\pi)$.
Away from the stationary point, write the integrand as
\[
 \frac1{\ii(\tau-2\pi z)}
 \frac{\mathrm d}{\mathrm dz}e^{\ii(\tau\log z-2\pi z)}.
\]
Integration by parts bounds an interval to the left of $L_0$ or to the right
of $R_0$ by $2/s$, since the reciprocal denominator is monotone there and
has modulus at most $1/s$ at the closest endpoint.
If the original interval lies wholly on one side, this already suffices.
Otherwise split it at $u=\max(a,L_0)$ and $v=\min(b,R_0)$. Each outside part
is bounded by $2/s$, with zero-length parts contributing zero. On $[u,v]$,
absolute integration gives
\[
 \frac{v-u}{a}\le\frac{R_0-L_0}{a}=\frac2s,
\]
using $s^2=2\pi a$. The three estimates sum to $6/s$.
This argument makes no assumption on the sign of $\tau$.
\end{proof}
\src{Oscillatory.lean; AllFrequency.lean}{norm_gmReflectionIntegral_le_left; norm_gmReflectionIntegral_le_right; norm_gmReflectionIntegral_le_six_div_sqrt}

For $m\in\mathcal B(L,H)$, change variables $z=mLx$ in $J_w(v,mL)$ and
apply Mellin inversion on the line $\Rea s=1$. All retained modes use the same
interval $[H/(40\pi),20H/\pi]$. Thus
\[
 J_w(v,mL)=\frac1{2\pi}\int_{\R}\Mell w(1+\ii r)
 (mL)^{\ii(r-v)}K(H,v,r)\dd r,
\]
where
\[
 K(H,v,r)=\int_{H/(40\pi)}^{20H/\pi}
 z^{-1}e^{\ii((v-r)\log z-2\pi z)}\dd z.
\]
Lemma~\ref{lem:allfreq} gives $|K(H,v,r)|^2\le2000/H$ for \emph{every}
$r$. Weighted Cauchy--Schwarz, with weight $|\Mell w(1+\ii r)|$, then
bounds the retained band by a squared short phase sum. Its Mellin mass is
finite by the compact-smooth Fourier decay. There is no discarded large-$r$
region and no hidden restriction $|r|\ll H$.

For a shell $H\le|t-u|<2H$, both signs are handled by conjugation of the real
cutoff. The retained integer band is partitioned into at most ten dyadic blocks.
The phase coefficients on each block have modulus one, and the complete-pair
Gram expansion dominates coefficients of modulus at most one by the
unweighted complete-pair energy. Localize $U$ to
\[
 U_{H,j}=U\cap[jH,(j+3)H),\qquad j\in\Z.
\]
Each point belongs to exactly three such sets; their diameters are at most
$3H$, hence their local diameters at most $4H$ when $H\ge1$.
Write $\mathscr R(L,H;U)$ for the exact finite energy
\[
 \mathscr R(L,H;U)=\sum_{M\in\mathcal B_{\rm dyad}(L,H)}
                     \sum_{j\in\mathcal J(U,H)}E_M(U_{H,j}),
\]
where $\mathcal B_{\rm dyad}(L,H)=\{2^{\lfloor\log_2m\rfloor}:m\in\mathcal B(L,H)\}$
and $\mathcal J(U,H)=\bigcup_{t\in U}\{\lfloor t/H\rfloor-2,\ldots,\lfloor t/H\rfloor\}$.
The latter is exactly the finite set of nonempty local-window indices.
The band block lengths range from $H/(40\pi L)$ to $4H/(\pi L)$;
the ratio is at most $160$, which places all their exponents in an interval
of at most ten integers. This proves the stated finite block bound.
The exact reflection estimate is
\begin{equation}\label{eq:reflection}
 \sum_{t,u\in U\atop H\le|t-u|<2H}|h_w(L,t-u)|^2
 \le\frac{20000}{\pi^2}\left(\int_{\R}|\Mell w(1+\ii r)|\dd r\right)^2
 \frac{L^2}{H}\mathscr R(L,H;U)
 +4C_A^2|U|^2L^{2-2A}.
\end{equation}
The diagonal has the separate bound $C_A^2|U|L^{2-2A}$.
\src{MellinKernel.lean; BandEnergy.lean; LocalCover.lean; Reflection.lean}{band_modeIntegral_mellin; norm_reflectionKernel_sq_le; norm_bandContribution_sq_le; sum_card_localSet; uniform_reflection}

\subsection{A finite supremum and uniform absorption}
For fixed nonempty admissible data define the finite-valued supremum
\begin{equation}\label{eq:bootstrapB}
 B=\sup_{\varnothing\ne U\subset W\atop N\le L\le8(T+1)}
 \frac{Q_L(U)}{b|U|(L^2+Nd_U)}.
\end{equation}
This is not an assumed kernel bound. The family is nonempty. Since $V\le N$
for nonempty $W$, $\Delta(T)\ge1$ and $L\ge N$, the denominator is at least
$1$. A crude finite-support bound on $h_w(L,v)$, uniform for
$0<L\le8(T+1)$, bounds each numerator by a fixed multiple of $|W|^2$.
Therefore the supremum is finite and nonnegative, and every descendant kernel
is bounded by its defining normalization times $B$.
\src{BootstrapDomain.lean; BootstrapSupremum.lean}{bootstrapFamily; crude_kernel_bound; bootstrapValues_bddAbove; bootstrapBound; kernelEnergy_le_bootstrapBound}

\begin{lemma}[Bootstrap inequality]\label{lem:boot}
There are constants $C_0,C_1\ge0$ depending only on $w$ such that
\[
 B\le C_0+C_1Nb^2\log(2T)B.
\]
For fixed $\kappa>0$, eventually uniformly in all admissible data,
$Nb^2\le1$ and $B\le2C_0$.
\end{lemma}
\begin{proof}
Apply the comparison lemma to every reflected dyadic block and every local set
in \eqref{eq:reflection}. Their new scale is still in
$[N,8(T+1)]$, so its kernel is controlled by the same finite supremum.
For a reflected dyadic length $M$, the band geometry gives
$(L^2/H)M^2\le16H$. The total local-cardinality sum is $3|U|$ and there are
at most ten reflected blocks. Consequently each shell is bounded by
\[
 CbNH|U|+CBb^3|U|(N^2H+NL^2)+E T^{-100}.
\]
The source obtains the error by taking a sufficiently high integration-by-parts
order and using $L\ge N\ge T^{1/3}$, $|U|\le T+1$.
Explicitly the off-band recursion uses $R_0=H+|k|$ and normalized
$p=(v/R_0,2\pi k/R_0)$ in the compact set
\[
 |p_1|\le2,\quad |p_2|\le2\pi,\quad
 |p_1-p_2x|\ge1/300\quad(1/3\le x\le6).
\]
Define $a_0=w$ and
$a_{j+1}(x)=\mathrm d[ a_j(x)x/(p_1-p_2x)]/\mathrm dx$.
All supports remain in $[1/2,5]$. Joint smoothness on this compact parameter
set bounds $a_j$ uniformly. If $I_j$ is its integral against
$e^{\ii R_0(p_1\log x-p_2x)}$, integration by parts gives
$R_0\ii I_j=-I_{j+1}$, with zero boundary terms.
Thus $|I_0|\le6C_j/R_0^j$. Summing its integer-mode majorants and choosing
$j$ large enough yields the stated shell error before any supremum absorption.
The one-separated pairs have no nonzero gap below $1$ and no gap exceeding
$\diam U$. Summing dyadic gap shells gives
$\sum H\le2d_U$ and at most $\log(2T)/\log2$ shells. With the diagonal added,
\[
 Q_L(U)\le2CbNd_U|U|
 +CBb^3|U|\left(2N^2d_U+NL^2\frac{\log(2T)}{\log2}\right)
 +E'T^{-90}.
\]
Divide by $b|U|(L^2+Nd_U)$, whose lower bound also absorbs the last error,
and take the supremum. The first normalized term is bounded by a constant;
the recursive term is at most a constant times $Nb^2\log(2T)B$.
This proves the first assertion without any circular assumption.

The height condition gives
\[
 Nb^2=\frac{N^3\Delta(T)^2}{V^4}
 \le\Delta(T)^2N^{-4\kappa}.
\]
Use $\Delta(T)\le D_\kappa T^{\kappa/6}$ and $N\ge T^{1/3}$ to obtain
$Nb^2\le D_\kappa^2T^{-\kappa}$. Both this expression and its product
with $\log(2T)$ tend to zero. A single threshold depending on $\kappa,w$
therefore makes $Nb^2\le1$ and $C_1Nb^2\log(2T)\le1/2$ for every
admissible configuration. Since $B\ge0$, the inequality implies
$B\le C_0+B/2$, hence $B\le2C_0$.
\end{proof}
\src{BootstrapShell.lean; BootstrapAbsorption.lean; UniformThreshold.lean}{bootstrap_one_shell; bootstrap_preabsorption; bootstrap_absorption_inequality; uniform_scalar_decay; uniform_absorbed_kernel}

\begin{proof}[Proof of Theorem~\ref{thm:LV}]
The absorbed kernel satisfies
$Q_L(W)\le Kb|W|(L^2+Nd_W)$ uniformly for $N\le L\le8(T+1)$.
Combine this at $L=NM$ with Lemma~\ref{lem:comparison} and
$E_M(W)\ge M|W|^2$. Divide by $M|W|$ for nonempty $W$ and use $Nb^2\le1$:
\[
 |W|\le(C_w+K)bNM+Kb^3\frac{Nd_W}{M}.
\]
Choose the positive integer $M=\lceil b\sqrt T\rceil$. Its rounding bounds
$b\sqrt T\le M\le b\sqrt T+1$ and $N\le T$ show that
$N\le NM\le8(T+1)$. The source explicitly derives
$Nb\sqrt T\le T$ from $Nb^2\le1$ and $N\le T$. Furthermore,
\[
 bNM\le Nb+Nb^2\sqrt T,
 \qquad b^3\frac{Nd_W}{M}\le2Nb^2\sqrt T,
\]
since $d_W\le T+1\le2T$. This gives, above the uniform threshold,
\[
 |W|\le A\left(\Delta(T)\frac{N^2}{V^2}
                   +\Delta(T)^2\frac{N^3\sqrt T}{V^4}\right).
\]
The divisor maximum estimate with loss $\varepsilon/2$ bounds both
$\Delta(T)$ and $\Delta(T)^2$ by $D^2T^\varepsilon$, with $D\ge1$.
For heights below the threshold, one-separatedness gives $|W|\le T+1$.
If $W\ne\varnothing$, $V\le N$, so the bracket
$N^2/V^2+N^3\sqrt T/V^4$ is at least one, and $T^\varepsilon\ge1$.
Adding the bounded-height threshold and $1$ to the large-height constant
therefore handles every $T\ge2$. The empty set has zero cardinality and
satisfies the same inequality immediately.
\end{proof}
\src{LargeValuesOptimization.lean; LargeValues.lean}{cardinality_at_length; optimal_integer_scale; large_values_with_divisors; large_values_data}

\section{The fourth moment from a genuine squared-zeta source}
\label{sec:fourth}
This part is not an assumed mean-value theorem. The source first constructs an
exact one-sided contour representation, then truncates it to an ordinary-divisor
polynomial and proves the polynomial's mean-square bound.

\subsection{Entire numerator and exact right-hand pieces}
Let $\Lambda_0$ denote mathlib's entire regular part of the completed zeta, so
that, away from $0,1$,
\[
 \Lambda(s)=\Lambda_0(s)-s^{-1}-(1-s)^{-1}.
\]
Define the entire function
\[
 \Xi_0(s)=s(1-s)\Lambda_0(s)-1.
\]
For $s\ne0,1$, it equals $s(1-s)\Lambda(s)$ and satisfies
$\Xi_0(1-s)=\Xi_0(s)$. In particular it is the pole-cleared numerator, not a
new independently specified zeta function. Put
\begin{gather}
 H_*(w)=(1-4w^2)^4,\qquad A_t=(s(t)(1-s(t)))^2=(1/4+t^2)^2,\\
 F_t(w)=\frac{e^{100w^2}H_*(w)\Xi_0(s(t)+w)^2}{A_t},\qquad
 \mathscr I_t(w)=F_t(w)/w.
\end{gather}
The denominator $A_t$ never vanishes. The numerator $F_t$ is entire,
$F_t(0)=\Lambda(s(t))^2$, and $F_t(-w)=F_{-t}(w)$.
\src{Stronger/Fourth/SquareSource.lean}{completedXiNumerator; hughesYoungAuxiliaryZero; zetaSquarePoleNormalization; zetaSquareContourNumerator; zetaSquareContourIntegrand}

\begin{lemma}[Exact source decomposition]\label{lem:squaresource}
For every $c>0$,
\[
 \Lambda(s(t))^2=J_c(t)+J_c(-t),\qquad
 J_c(t)=\frac1{2\pi}\int_{\R}\mathscr I_t(c+\ii u)\dd u.
\]
Writing $B_t=\Gamma_{\R}(s(t))\Gamma_{\R}(s(-t))$, define
\[
 R(t)=J_1(-t)/B_t.
\]
Then $B_t\ne0$, $\zabs(t)^2=2\Rea R(t)$, and
$\zabs(t)^4\le4|R(t)|^2$.
\end{lemma}
\begin{proof}
Integrate $\mathscr I_t$ around the rectangle with vertical edges $-c,c$ and
horizontal heights $\pm U$. The only pole is the simple pole at $0$ and its
residue is $F_t(0)$. Reflection $w\mapsto-w$ converts the left vertical
edge to the right vertical edge for $-t$. The source proves a polynomial
bound for $\Xi_0$ on each fixed vertical strip and a Gaussian-polynomial
majorant for $F_t(x+\ii u)$, uniform for $|x|\le c$. It follows that the
horizontal edges tend to zero and both vertical integrals are absolutely
integrable. Taking the limit of the finite rectangle identity gives the first
formula, with its normalization $1/(2\pi)$ (equivalently $1/(2\pi\ii)$
against $\mathrm dw$).

On the critical line the functional equation gives
$\Lambda(s(t))=\Lambda(s(-t))$, whereas conjugation gives
$\Lambda(s(-t))=\overline{\Lambda(s(t))}$. Thus the completed critical-line
value is real. The Gamma factors are nonzero, and
\[
 \Lambda(s(t))^2=\zabs(t)^2B_t.
\]
Conjugation of the right-hand integrands gives $J_1(t)=\overline{J_1(-t)}$
and $B_t\in\R$. Dividing the exact source formula by $B_t$ therefore gives
$\zabs(t)^2=2\Rea R(t)$. Squaring and using $|\Rea z|\le|z|$ proves the
last inequality.
\end{proof}
\src{Stronger/Fourth/SquareSource.lean; SquareContour.lean; SquareNormalized.lean}{zetaSquare_finiteRectangle; completedZeta_square_eq_vertical_integrals; zetaSquareNorm_eq_two_re_fourthRightPiece; zeta_fourth_le_four_mul_rightPiece_sq}

The arithmetic identity used to open the right contour is
\[
 \zeta(s)^2=\sum_{n\ge1}d(n)n^{-s},\qquad \Rea s>1.
\]
It is proved locally from the Dirichlet convolution of the constant-one
arithmetic function with itself, which counts divisors, and multiplication of
absolutely convergent Dirichlet series. It is not inferred at $\Rea s=1/2$.
For a shift $w$, put
\begin{align}
 P_t(w)&=\left(\frac{(s(-t)+w)(1-s(-t)-w)}{s(-t)(1-s(-t))}\right)^2,\\
 \mathcal K_t(w)&=\frac{e^{100w^2}H_*(w)}w P_t(w)
             \frac{\Gamma_{\R}(s(-t)+w)^2}{B_t}.
\end{align}
On $\Rea w=1$ the squared-zeta series is absolutely convergent and the exact
normalized source is
\[
 R(t)=\frac1{2\pi}\int_{\R}\mathcal K_t(1+\ii u)
              \sum_{n\ge1}d(n)n^{-s(-t)-1-\ii u}\dd u.
\]
The source proves summability of the integrals of term norms before interchanging
the sum and this integral.
\src{Stronger/Fourth/SquareSource.lean; SquareNormalized.lean; SquareDivisorSeries.lean}{riemannZeta_sq_eq_divisorLSeries; zetaFourthKernel; zetaFourthRightPiece_eq_divisor_integral; hasSum_zetaSquareDivisorContribution}

\subsection{Gamma shifts, finite prefixes and the tail}
The kernel estimates are derived, not supplied as an asymptotic shortcut. The
local digamma development starts from Euler's Gamma limit and the convergent
series
\[
 \psi_\Gamma(z)+\gamma=\sum_{m\ge0}\left(\frac1{m+1}-\frac1{z+m}\right),
 \qquad \Rea z>0,
\]
where $\psi_\Gamma=\Gamma'/\Gamma$ and $\gamma$ is Euler's constant. Comparing
the reciprocal sum with the telescoping logarithmic integral gives
\[
 |\psi_\Gamma(z)-\log z|\le4/|\Ima z|\qquad(|\Ima z|\ge1,\ \Rea z>0).
\]
The normalized shift
\[
 \left(\frac{\Gamma_{\R}(s(-t)+vw)}{\Gamma_{\R}(s(-t))}\right)^2
 e^{-vw\ell_t},\qquad
 \ell_t=\log(t/(2\pi))-\ii\pi/2,
\]
has initial value one and logarithmic derivative bounded by
$(17|w|+2|w|^2)/t$ when $t\ge4$, $\Rea w\ge0$, $|w|\le t/2$,
$0\le v\le1$. Integrating the derivative bound gives both its exponential
bound and its difference-from-one bound. The pole factor $P_t(w)$ is bounded
and has difference from one at most $12|w|/t$ on this range. For the far
$|u|$ range, reciprocal Gamma reflection and elementary exponential estimates
are multiplied by $e^{-100u^2}$. These give, for every fixed $c>0$, a constant
$K_c>0$ such that
\begin{equation}\label{eq:fourthkernel}
 |\mathcal K_t(c+\ii u)|\le K_c t^c e^{-98u^2}
 \qquad(t\ge4,\ t\ge4c,\ u\in\R).
\end{equation}
The estimates explicitly separate near and far shifts and prove their
integrability.
\src{Stronger/Fourth/DigammaSeries.lean; DigammaLog.lean; GammaShiftAmplitude.lean; SquareKernel.lean}{hasSum_digamma_of_re_pos; norm_digamma_sub_log_le; norm_gammaReal_shift_sq_le_exp; exists_norm_zetaFourthKernel_le}

For a finite set $S\subset\N$ and $c>0$, define
\[
 R_S(t;c)=\frac1{2\pi}\int_{\R}\mathcal K_t(c+\ii u)
       \sum_{n\in S}d(n)n^{-1/2-c+\ii(t-u)}\dd u.
\]
The $n=0$ term is defined to be zero. Each summand has no pole between two
positive shift lines; the Gamma arguments remain in the positive half-plane,
the explicit $1/w$ pole stays to the left, and the Gaussian kills horizontal
edges. Thus its integral is independent of the chosen positive line. Shift the
finite prefix to a small line $c$, and shift its complement to a large line
$b>3/2$. The exact identity is
\[
 R(t)=R_{\{0,\ldots,N\}}(t;c)+R_{\{n>N\}}(t;b).
\]
This is a decomposition of the actual source, not an approximate functional
equation assumed for the prefix. The complement is absolutely convergent on
$\Rea w=b$. Since $d(n)\le n$, comparison with an integral gives
\[
 \sum_{n>N}d(n)n^{-1/2-b}
 \le\frac{N^{3/2-b}}{b-3/2}.
\]
From \eqref{eq:fourthkernel},
\[
 |R_{\{n>N\}}(t;b)|\le K_b' t^bN^{3/2-b}.
\]
If $N\ge H^{1+q}$, $H\le t\le2H$ and $q>0$, choose
$b=2+2/q$. The identity
\[
 b+(1+q)(3/2-b)=-(1+q)/2
\]
then bounds the entire tail by a constant, uniformly in $H,t,N$ and in the
positive prefix shift $c$.
\src{Stronger/Fourth/SquareTerm.lean; SquareLineShift.lean; SquareTruncation.lean; SquareTail.lean}{zetaFourthRightPiece_eq_prefix_add_tail; tsum_nat_rpow_tail_le; exists_norm_fourthRightPiece_sub_cutoff_le}

\subsection{The polynomial mean square is proved explicitly}
\begin{lemma}[Dyadic Dirichlet mean square]\label{lem:meanvalue}
For arbitrary complex coefficients $a_n$, positive integer $N$ and $H\ge0$,
\[
 \int_0^H\left|\sum_{N<n\le2N}a_n n^{-\ii t}\right|^2\dd t
 \le\bigl(H+2(5\pi+1)N\bigr)\sum_{N<n\le2N}|a_n|^2.
\]
The same bound is invariant under translation or reflection of the integration
variable, with the interval length replacing $H$.
\end{lemma}
\begin{proof}
Write
\[
 Q_N(a)=\sum_{N<m,n\le2N\atop m\ne n}
 \frac{\overline{a_n}a_m}{\log m-\log n},\qquad
 a_n^{(H)}=a_n n^{-\ii H}.
\]
Integration of each exponential term gives the exact identity
\[
 \int_0^H\left|\sum a_n n^{-\ii t}\right|^2\dd t
 =H\sum|a_n|^2+\ii\bigl(Q_N(a^{(H)})-Q_N(a)\bigr),
\]
where the real left side is viewed inside $\C$.
For $x,y>0$, $x\ne y$, elementary bounds for $\log(y/x)$ give
\[
 \left|\frac1{\log y-\log x}-\frac{x}{y-x}\right|\le1.
\]
Consequently $Q_N$ is the sum of a main discrete Hilbert form and an error
whose norm is at most $(\sum|a_n|)^2\le N\sum|a_n|^2$.

For completeness, the discrete Hilbert form here is obtained directly from
Fourier orthogonality. For finite integer frequencies $S$ and coefficient
families $a,b$, define
\[
 \mathcal H_S(a,b)=\sum_{m,n\in S\atop m\ne n}
          \frac{\overline{a_n}b_m}{m-n},\quad
 A(x)=\sum_{n\in S}a_ne^{\ii nx},\quad B(x)=\sum_{n\in S}b_ne^{\ii nx}.
\]
Termwise integration by parts yields
\[
 \int_0^{2\pi}(\pi-x)\overline{A(x)}B(x)\dd x
     =2\pi\ii\mathcal H_S(a,b).
\]
Use $|\pi-x|\le\pi$, the elementary product bound
$|A||B|\le|A|^2+|B|^2$, and Fourier orthogonality
$\int_0^{2\pi}|A|^2=2\pi\sum|a_n|^2$ to obtain
\[
 |\mathcal H_S(a,b)|\le\pi\left(\sum|a_n|^2+\sum|b_n|^2\right).
\]
Take the scaled first family $\widetilde a_n=(n/N)a_n$, whose square mass is at most
$4\sum|a_n|^2$. The main logarithmic form is $N\mathcal H_S(\widetilde a,a)$,
so it has norm at most $5\pi N\sum|a_n|^2$. Adding the reciprocal-log error
gives $|Q_N(a)|\le(5\pi+1)N\sum|a_n|^2$. The endpoint twist preserves
square mass, so the two endpoint forms give the claimed factor $2$.
A translation only twists coefficients by unit phases; reflection conjugates
or reverses those phases. Thus the same estimate holds on translated intervals.
\end{proof}
\src{Stronger/Fourth/MeanValueCore.lean; PrefixMean.lean}{hilbertForm_norm_le; abs_inv_log_sub_sub_div_le_one; logHilbertQuad_norm_le; integral_norm_sq_dirichletTime_le; integral_norm_sq_dirichletPrefix_reflected_le}

\begin{proposition}[Native fourth moment]\label{prop:fourth}
For every $\eta>0$ there are $C\ge0$ and $H_0$ such that for $H\ge H_0$,
$H>0$,
\[
 \int_H^{2H}\zabs(t)^4\dd t\le C H^{1+\eta}.
\]
\end{proposition}
\begin{proof}
The source chooses $q=\min(\eta,1)/20$. The unparenthesized Lean expression
\lean{min eta 1/20} has this meaning because function application binds before
division; it must not be misread as $\min(\eta,1/20)$. Thus $0<q\le1/20$
and $7q\le\eta$.
Choose a dyadic length $N=2^M$ with
$H^{1+q}\le2^M\le2H^{1+q}$. The preceding tail bound gives
$|R(t)-R_{\{0,\ldots,N\}}(t;q)|\le K_t$ uniformly on $[H,2H]$.
Define
\[
 a_n=d(n)n^{-1/2-q},\qquad
 F(t)=\int_{\R}e^{-98u^2}\left|\sum_{1\le n\le N}a_n n^{-\ii(u-t)}\right|^2\dd u.
\]
The finite prefix source and \eqref{eq:fourthkernel}, followed by weighted
Cauchy--Schwarz with Gaussian mass $\sqrt{\pi/98}$, imply
$|R_{\{0,\ldots,N\}}(t;q)|^2\le K_g t^{2q}F(t)$.
Lemma~\ref{lem:squaresource} and $|x+y|^2\le2|x|^2+2|y|^2$ give
\[
 \zabs(t)^4\le8K_g(2H)^{2q}F(t)+8K_t^2.
\]
All Gaussian-product functions are integrable: a finite polynomial is bounded
by its coefficient norm sum, and the Gaussian is integrable. Fubini may
therefore be applied.

The divisor subpower estimate gives, for every $q>0$, the uniform bound
\[
 \sum_{L<n\le2L}|d(n)n^{-1/2-c}|^2\le D(2L)^q\qquad(c\ge0).
\]
Partition the prefix into its $n=1$ term and dyadic blocks. Finite
Cauchy--Schwarz costs $(M+1)^2$ after the common block bound is used.
Lemma~\ref{lem:meanvalue} then gives
\[
 \int_H^{2H}F(t)\dd t
 \le D_p(M+1)^2\bigl(H+2(5\pi+1)2^M\bigr)(2^M)^q.
\]
The source absorbs $(M+1)^2$ into a fixed multiple of $(2^M)^q$.
With $2^M\le2H^{1+q}$ the full scaled expression is bounded by
$D_b H^{1+7q}$: before the final simplification its power is
$1+5q+2q^2\le1+7q$ for $q\le1$.
Integrating the pointwise inequality gives
\[
 \int_H^{2H}\zabs(t)^4\dd t
 \le8K_gD_pD_bH^{1+7q}+8K_t^2H
 \le(8K_gD_pD_b+8K_t^2)H^{1+\eta}.
\]
The height threshold is chosen so that every kernel, tail and dyadic-cutoff
condition holds before $t$ varies. This is the native producer consumed later.
\end{proof}
\src{Stronger/Fourth/GaussianProducer.lean; SquarePrefix.lean; FourthMoment.lean}{exists_integral_gaussian_fourthPolynomial_le; exists_norm_sq_fourthPrefix_le_gaussian; exists_fourth_dyadic_budget; zeta_fourth_dyadic}

\section{Point values controlled by actual weighted local second means}
\label{sec:pointmean}
For $t\in\R$ and $L\ge0$, define
\[
 M(t,L)=\int_{-L}^{L}e^{-|u|}\zabs(t+u)^2\dd u
       =\int_{t-L}^{t+L}e^{-|t-v|}\zabs(v)^2\dd v.
\]
The second equality is exactly translation of the first, with no kernel
replacement. Also put
\[
 M_\infty(t)=\int_{\R}e^{-|u|}\zabs(t+u)^2\dd u,\quad
 M_\delta^*(t)=\int_{\R}\frac{e^{-4|u|/3}}{\delta+|u|}\zabs(t+u)^2\dd u
 \quad(\delta>0).
\]
All these integrals are genuine convergent integrals by coarse quadratic growth
and exponential decay.

\begin{proposition}[Native point-to-mean inequality]\label{prop:pointmean}
There is $C>0$ such that, for every real $t\ge10$,
\[
 \zabs(t)^2\le C\log t\bigl(1+M(t,(\log t)^2)\bigr).
\]
\end{proposition}
\begin{proof}
Here the source parameter is $\delta=1/\log t$ and the rectangle height is
$R=(\log t)^2$. The initial finite rectangle integrates
$\Gamma(w)\zeta(s(t)+w)^2$ over $-\delta\le\Rea w\le\delta$,
$|\Ima w|\le R$. Only the Gamma pole at zero is enclosed; its residue is
$\zeta(s(t))^2$.

To bound its vertical sides, the source first establishes an off-critical
inequality. For $s=1/2+\delta+\ii t$, Mellin inversion of the actual
exponentially smoothed divisor series gives
\[
 \zeta(s)^2=\sum_{n\ge1}d(n)n^{-s}e^{-n}
     -\operatorname{Res}_{w=1-s}\bigl(\Gamma(w)\zeta(s+w)^2\bigr)
     -\frac1{2\pi}\int_{\R}\Gamma(-\delta+\ii u)
                    \zeta(1/2+\ii(t+u))^2\dd u.
\]
The shift crosses the Gamma pole at zero and the moving double zeta pole
$w=1-s$. The latter is evaluated by differentiating its cleared numerator:
if $\zeta_1(z)=z\zeta(1+z)$ is the analytic pole-removed function with
$\zeta_1(0)=1$, its residue is
\[
 \Gamma(1-s)\psi_\Gamma(1-s)+2\Gamma(1-s)\zeta_1'(0).
\]
The smoothed series is bounded absolutely by $\sum n e^{-n}$; the residue is
uniformly bounded for $0<\delta\le1/4$, $t\ge10$ using Gamma decay and a
digamma bound. All horizontal limits and vertical integrability are proved
with exponential or polynomial strip majorants. Hence
\[
 |\zeta(1/2+\delta+\ii t)|^2\le C\bigl(1+M_\delta^*(t)\bigr).
\]
The functional equation transfers this to $1/2-\delta+\ii t$, with an
additional factor at most
$\exp(2\delta(\log(t/2+2)+D))$. This factor is obtained by integrating a
Gamma logarithmic-derivative bound on the fixed positive-real-part strip.
It is not assumed to be one.

The Gamma bounds on the small vertical lines are
\[
 (\delta+|u|)|\Gamma(\pm\delta+\ii u)|\le144e^{-4|u|/3}.
\]
After substitution of the off-critical inequalities, the vertical-edge
majorant contains a convolution of two kernels
$k_\delta(u)=e^{-4|u|/3}/(\delta+|u|)$.
The source explicitly proves
\[
 \int_{\R}k_\delta(w)k_\delta(u-w)\dd w
 \le\frac\pi\delta e^{-|u|}.
\]
Indeed the excess exponential reserve bounds the product by
$e^{-|u|}/2$ times the sum of the squares of
$r_\delta(x)=e^{-|x|/6}/(\delta+|x|)$ at $w$ and $u-w$.
Each square is at most $(\delta^2+x^2)^{-1}$, whose integral is
$\pi/\delta$. Tonelli and the proved product majorants then yield the
convolution inequality and bound both vertical edges by
$C\delta^{-1}(1+M_\infty(t))$.

Eventually $R\le t/2$, so every shifted height on the edges is still large.
For $|u|\le R$ the reflected Gamma factor above is at most $e^4$ when
$D\le\log t$; this is proved before it is absorbed into the constant.
The horizontal edges have bounds
$C e^{-4R/3}(5(1+t+R))^2$ and are eventually at most one, since
$R=(\log t)^2$. The residue rectangle therefore gives
\[
 \zabs(t)^2\le C\log t(1+M_\infty(t))
\]
eventually.

Finally, if $|u|\ge(\log t)^2$ and $t$ is sufficiently large, coarse quadratic
growth gives
$e^{-|u|}\zabs(t+u)^2\le19C_0e^{-|u|/2}$.
Its tail integral is uniformly bounded. Thus
$M_\infty(t)\le M(t,(\log t)^2)+D_0$, and the last additive constant is
absorbed into $1+M$. Continuity of $\zabs$ bounds the remaining compact range
$10\le t\le t_0$; since $\log t\ge\log10>0$ and $M\ge0$, increasing $C$
extends the estimate to every $t\ge10$.
\end{proof}
\src{Stronger/PointMean/FiniteRectangle.lean; MellinShift.lean; DoubleConvolution.lean; VerticalEdges.lean; Native.lean}{heathBrownLemmaThree_finiteRectangle; heathBrown_twoPole_vertical_shift; integral_strongKernel_mul_strongMoment_le; exists_eventually_heathBrownLemmaThree_fullMoment; heathBrownLemmaThree_native}

\begin{lemma}[Exponential overlap and the finite-peak inequality]\label{lem:equation44}
For a one-separated finite $W$ and any $x\in\R$,
$\sum_{t\in W}e^{-|t-x|}\le4$.
Suppose $W\subset[c-G/2,c+G/2]$, $G,L\ge0$, $G/2+L\le G$,
$c-G/2\ge10$, $c+G/2\le T$, $T\ge10$,
$(\log t)^2\le L$ and $\zabs(t)\ge V>0$ for $t\in W$. Then
\[
 V^2|W|\le C\log T\left(|W|+4\int_{c-G}^{c+G}\zabs(u)^2\dd u\right).
\]
\end{lemma}
\begin{proof}
For each integer $k\ge0$, at most two points of $W$ have
$k\le|t-x|<k+1$, one on each side, by one-separatedness. On this shell
$e^{-|t-x|}\le2^{-k}$. Summing $2\sum_{k\ge0}2^{-k}=4$ proves the overlap
bound. Proposition~\ref{prop:pointmean} gives the displayed point inequality
at each $t$; enlarge its integration radius from $(\log t)^2$ to $L$ and
bound $\log t\le\log T$. Every resulting interval $[t-L,t+L]$ lies in
$[c-G,c+G]$ by the fit condition. Interchange the finite sum and integral,
apply the overlap bound to the nonnegative integrand, and sum the constant
terms. This is exactly the source's equation-(44) packet interface.
\end{proof}
\src{Stronger/PointMean/ExponentialOverlap.lean; IntegralOverlap.lean; Equation44.lean}{sum_exp_neg_abs_sub_le_four; sum_truncated_exp_kernel_zeta_sq_le_localSecondMoment; heathBrown_equation44_native}

\section{The actual Atkinson source and both stationary contributions}
\label{sec:atkinson}
The second analytic producer needed for the high-value bound is an estimate for
packets of local second means. Its proof retains an exact compact-test Voronoi
identity, both Bessel branches, an actual quadratic Gaussian and both signed
stationary carriers. The construction is recorded here because replacing it by
an unspecified ``Atkinson formula'' would conceal substantive formal content.

\subsection{Frozen Gamma source and physical Gaussian}
Put
\[
 \chi_t=\frac{\Gamma_{\R}(s(-t))}{\Gamma_{\R}(s(t))},\qquad |\chi_t|=1,
 \qquad\ell_t=\log(t/(2\pi))-\ii\pi/2.
\]
The divisor weight is the entire function
\begin{equation}\label{eq:divisorweight}
 W(q)=\frac1{2\pi}\int_{\R}
       \frac{e^{100(1+\ii u)^2}H_*(1+\ii u)e^{-q(1+\ii u)}}{1+\ii u}\dd u.
\end{equation}
It obeys $W(q)+W(-q)=1$. To see this, shift the line from $1$ to $-1$;
the numerator is entire, its value at zero is one, and the only pole of its
quotient by $w$ is at zero. Gaussian horizontal decay permits the limit of the
finite rectangle. The left line, after $w\mapsto-w$, is the negative of the
right line for $-q$. Thus their difference is one. The same Gaussian majorants
justify differentiation under the integral, so $W$ is entire, with
$W'(q)=-(2\pi)^{-1}\int e^{100w^2}H_*(w)e^{-qw}\dd u$.
On a fixed imaginary strip it satisfies
\begin{gather}
 |W(q)|\le C\min(1,e^{-\Rea q}),\\
 |W(q+h)-W(q)|\le C|h|e^{|h|}\min(e^{-\Rea q},e^{\Rea q})
 \quad(h\in\R).
\end{gather}
The first decay follows on the right contour and then by reflection;
the second follows by integrating the derivative and using the same
reflection. These estimates precede all physical heights and indices.
\src{Stronger/Atkinson/DivisorWeightKernel.lean; DivisorWeightReflection.lean; DivisorWeightVariation.lean}{zetaDivisorWeight; zetaDivisorWeight_add_neg; exists_norm_zetaDivisorWeight_min_le; exists_norm_zetaDivisorWeight_real_sub_min_le}

The Gamma shift calculation from Section~\ref{sec:fourth} approximates the
normalized right source by
\[
 \chi_t\sum_{n\ge1}d(n)n^{-1/2+\ii t}W(\log n-\ell_t)
\]
with a bounded error. The exact near-kernel error has an integrable majorant
$C e^{-90u^2}(1+|u|)^{12}$, and the far-kernel part has the same form.
Coefficientwise errors sum because $\sum d(n)n^{-3/2}<\infty$.
Freeze the weight at a center $T$ to define
\[
 F(T,x)=\chi_{T+x}\sum_{n\ge1}d(n)n^{-1/2+\ii(T+x)}W(\log n-\ell_T).
\]
For $T\ge8$, $|x|\le T/2$, every fixed $\theta>0$ gives
\begin{equation}\label{eq:frozenerror}
 |\zabs(T+x)^2-2\Rea F(T,x)|\le C_\theta(1+|x|T^{-1/2+\theta}).
\end{equation}
Indeed $\log(T+x)-\log T$ has absolute value at most $2|x|/T$,
and the weight difference bound can be summed against
$d(n)n^{-1/2}\min(1,T/(2\pi n))$; this mass is
$O_\theta(T^{1/2+\theta})$. This is the source's freezing error, not a
formal replacement of the height-dependent phase by a constant.
\src{Stronger/Atkinson/KernelApproximation.lean; DivisorWeightSource.lean; DivisorWeightFreezing.lean; RealFrozenSource.lean}{exists_integral_norm_zetaSquareRightKernel_error_le; zetaSquareLeadingDivisorContribution_eq_weight; exists_abs_zetaSquareNorm_sub_frozen_le}

For $G>0$ define
\[
 A(T,G)=G^{-2}+\ii/(2T),\quad
 Q(T,G;v)=\int_{\R}e^{\ii vx}e^{-A(T,G)x^2}\dd x
       =\left(\frac\pi{A(T,G)}\right)^{1/2}e^{-v^2/(4A(T,G))}.
\]
The real part of $A$ is positive; the principal half-power is therefore the
one in the convergent complex Gaussian formula. If $G^2\le2T$, then
\begin{equation}\label{eq:Gaussianbound}
 |Q(T,G;v)|\le\sqrt\pi G e^{-(Gv)^2/8}.
\end{equation}
This follows from $|\Ima A|\le\Rea A$,
$\Rea A^{-1}\ge1/(2\Rea A)$ and the norm of the half-power.

The logarithmic derivative of $\chi_t$ is
$\ii(\log\pi-\Rea\psi_\Gamma(s(t)/2))\chi_t$, and the quantity in parentheses
is $-\log(t/(2\pi))+O(1/t)$. On $|x|\le r\le T/2$ the source proves
\[
 |\chi_{T+x}-\chi_Te^{-\ii(x\log(T/(2\pi))+x^2/(2T))}|
 \le(18/T+2r^2/T^2)|x|.
\]
Integrating this error on $[-r,r]$ against the physical Gaussian
$e^{-(x/G)^2}$ costs at most $2r^2(18/T+2r^2/T^2)$.
The two omitted Gaussian tails cost
$2\sqrt{2\pi}G e^{-(r/G)^2/2}$. Thus the actual phase-Gaussian transform
is compared to $\chi_T Q(T,G;v-\log(T/(2\pi)))$, with both errors stated
and proved. Summability of frozen coefficients and product Gaussian
majorants justify the integral/series exchanges.
\src{Stronger/Atkinson/QuadraticGaussian.lean; GammaQuadratic.lean; GammaGaussianTransform.lean; FrozenDivisorGaussian.lean}{zetaGaussianQuadraticIntegral_eq; norm_zetaGaussianQuadraticIntegral_le; norm_zetaSquareGammaGaussianTransform_sub_quadratic_le; hasSum_zetaFrozenDivisorGaussianMean}

\subsection{The compact source test and the lattice phase}
Let $S_0$ be mathlib's smooth transition: it is zero on $(-\infty,0]$,
one on $[1,\infty)$ and takes values in $[0,1]$. For $L>0$ set
\[
 b_v=\frac{T}{2\pi}e^{v/G},\qquad
 \vartheta_{T,G,L}(x)=S_0\!\left(\frac{x-b_{-2L}}{b_{-L}-b_{-2L}}\right)
 S_0\!\left(\frac{b_{2L}-x}{b_{2L}-b_L}\right).
\]
It equals one on $[b_{-L},b_L]$ and is supported in $[b_{-2L},b_{2L}]$.
The source's compact test is
\begin{equation}\label{eq:compacttest}
 g_{T,G,L}(x)=\vartheta_{T,G,L}(x)x^{-1/2+\ii T}
 W(\log x-\ell_T)\chi_TQ(T,G;\log x-\log(T/(2\pi)))e^{-2\pi\ii x}
 \quad(x>0),
\end{equation}
and is zero off the positive support interval. The phase $e^{-2\pi\ii x}$
is not cosmetic: it is identically one on the integer lattice, so it preserves
the discrete divisor sum, but it is retained in every subsequent continuous
Voronoi integral. Removing it would alter the saddles.
Define $\Sigma(T,G,L)=\sum_{n\ge1}d(n)g_{T,G,L}(n)$.
The source constructs $g$ as a \lean{DFIVoronoiTestFunction}, consisting of
positive lower support endpoint, upper endpoint, their ordering,
$C^\infty$ smoothness and support inclusion. Hence its Mellin transform is
entire and decreases faster than any fixed power on each vertical line.
\src{Stronger/Atkinson/DivisorBandCutoff.lean; SmoothDivisorTest.lean; DivisorLatticePhase.lean; AtkinsonAmplitude.lean}{zetaDivisorBandCutoff; zetaSmoothDivisorVoronoiTest; zetaAtkinsonDivisorTest; zetaAtkinsonDivisorTest_eq_amplitude_phase}

For every fixed $\delta>0$, eventually for
$T^\delta\le G\le T^{1/2-\delta}$, the preceding Gaussian source construction,
with $L=\log T$, yields
\begin{equation}\label{eq:physicalsource}
 \left|\int_{\R}e^{-((t-T)/G)^2}\zabs(t)^2\dd t
                  -2\Rea\Sigma(T,G,\log T)\right|\le C_\delta G\log T.
\end{equation}
The source obtains this by integrating \eqref{eq:frozenerror}, replacing the
actual Gamma Gaussian transform by its quadratic one, restricting to
$|G(\log n-\log(T/(2\pi)))|\le\log T$, and smoothing the band. Each omitted
term is bounded before summation; \eqref{eq:Gaussianbound} makes the band
complement exponentially small. The freezing and quadratic errors use fixed
positive losses smaller than the width margin $\delta$.
On $[T-G,T+G]$, $e^{-((t-T)/G)^2}\ge e^{-1}$, so
\eqref{eq:physicalsource} also implies
\begin{equation}\label{eq:localsource}
 I(T,G):=\int_{T-G}^{T+G}\zabs(t)^2\dd t
          \le2e\Rea\Sigma(T,G,\log T)+C_\delta G\log T.
\end{equation}
\src{Stronger/Atkinson/GaussianAveraging.lean; SmoothDivisorSource.lean}{exists_zetaSquareGaussianMean_smooth_approximation; exists_zetaSquareLocalMean_le_smooth_divisor}

\subsection{Native modulus-one Voronoi, including both complete dual series}
\begin{lemma}[Actual compact-test divisor Voronoi]\label{lem:voronoi}
Let $g:\R\to\C$ be $C^\infty$ and supported in $[a,b]$ with $0<a\le b$.
Then
\begin{align}\label{eq:voronoi}
 \sum_{n\ge1}d(n)g(n)
 &=\int_0^\infty(\log x+2\gamma)g(x)\dd x\\
 &\quad-2\pi\sum_{n\ge1}d(n)\int_0^\infty
                          g(x)Y_0(4\pi\sqrt{nx})\dd x\nonumber\\
 &\quad+4\sum_{n\ge1}d(n)\int_0^\infty
                          g(x)K_0(4\pi\sqrt{nx})\dd x.\nonumber
\end{align}
Both dual series are convergent series with the displayed terms.
\end{lemma}
\begin{proof}
The initial Mellin integral is
$(2\pi\ii)^{-1}\int_{(2)}\zeta(s)^2\Mell g(s)\mathrm ds$.
Termwise Mellin inversion gives the divisor sum, since the coefficient series
is absolutely convergent on this line. Shift to $\Rea s=-1/2$.
The double pole at $s=1$ has cleared numerator
$((s-1)\zeta(s))^2\Mell g(s)$. Its derivative at $1$ equals
$\Mell g'(1)+2\gamma\Mell g(1)$, hence the integral main term in
\eqref{eq:voronoi}. Smooth compact support proves arbitrary Mellin vertical
decay and entireness, and coarse zeta strip growth makes the horizontal edges
vanish; the infinite-line shift is therefore a limit of proved finite
rectangle identities.

The actual functional equation splits the reflected squared-zeta multiplier
into two exact multipliers. If
$A(s)=(2\pi)^{-s}\Gamma(s)$, these are
\[
 M_-(z)=A(1-z)^2(e^{\pi\ii(1-z)}+e^{-\pi\ii(1-z)}),\qquad
 M_+(z)=2A(1-z)^2.
\]
For $n\ge1$ the corresponding dual terms are
\[
 \mathcal V_\pm g(n)=\frac1{2\pi\ii}\int_{(-1/2)}
           n^{z-1}M_\pm(z)\Mell g(z)\mathrm dz.
\]
Both multipliers are holomorphic on $\Rea z<1$; on the initial dual line their
norms are bounded by a quadratic polynomial in $|\Ima z|$. Absolute Mellin
integrability, followed by $\sum d(n)n^{-3/2}<\infty$, permits both complete
series to be opened and added separately.

The source defines the kernels by actual integrals,
\begin{align*}
 K_0(x)&=\int_0^\infty e^{-x\cosh u}\dd u,\\
 Y_0(x)&=\frac2\pi\left(\int_0^{\pi/2}\sin(x\cos u)\dd u
                -\int_0^\infty\frac{e^{-xu}}{\sqrt{1+u^2}}\dd u\right),\quad x>0.
\end{align*}
Gaussian domination establishes the first integral's convergence; exponential
domination establishes the second tail's convergence, and the oscillatory
integral is over a finite interval.
Fubini and the beta integral give
$\Mell K_0(s)=2^{s-2}\Gamma(s/2)^2$ for $\Rea s>0$.
For $Y_0$, the source proves its Mellin identity first in the absolute strip
$1/4<\Rea s<1/2$ and identifies the continued symbol by
\[
 \Mell Y_0(2w)=-\frac2\pi\cos(\pi w)\Mell K_0(2w)
\]
where the left-hand Mellin integral is used only in its convergence strip.
To pair the actual $Y_0$ kernel with $g$, move the dual line to
$\Rea z=13/16$, so $2(1-z)$ has real part $3/8$ inside that strip.
Mellin pairing then gives $\mathcal V_-g(n)=-2\pi\int g(x)Y_0(4\pi\sqrt{nx})\dd x$.
Moving back is justified by holomorphy and proved decay; it is not an assumed
identification of a divergent Mellin integral. The analogous $K_0$ pairing
gives $\mathcal V_+g(n)=4\int g(x)K_0(4\pi\sqrt{nx})\dd x$.
Substitution into the two complete dual series proves \eqref{eq:voronoi}.
\end{proof}
\src{Stronger/Atkinson/OrdinaryDivisorVoronoi.lean; BesselK0Mellin.lean; BesselY0Mellin.lean; BesselTransformBridge.lean; VoronoiMinusTransport.lean; ZetaAtkinsonVoronoi.lean}{ordinaryDivisorVoronoi_native; dfiVoronoiPlusTransform_mellin_eq_bessel; verticalIntegral_dfiVoronoiMinus_mellin_thirteenSixteenths_eq_bessel; ordinaryDivisorVoronoi_bessel}

\paragraph{What is removed, and why.}
Apply Lemma~\ref{lem:voronoi} to \eqref{eq:compacttest}. The logarithmic main
term is retained initially. The $K_0$ branch is bounded as a complete series
using its exponential large-argument decay and the fact that its positive
support is comparable to $T$. The source proves a power saving
$O_{\delta,A}(GT^{-A})$ for every fixed $A$; in the native local-mean
composition only the $A=0$ consequence is needed. The main logarithmic
integral is bounded by $O_\delta(G\log T)$ using its exact carrier phase and
its smooth amplitude, not by replacing the divisor sum with an assumed
asymptotic main term. Both bounds are absorbed into the existing error in
\eqref{eq:localsource}.
\src{Stronger/Atkinson/AtkinsonK0Scales.lean; AtkinsonMainIntegral.lean; ZetaAtkinsonMinusSource.lean}{exists_zetaAtkinsonBesselPlus_powerSaving; exists_zetaSquareLocalMean_le_atkinson_minus}

\subsection{Two-term Neumann expansion and actual saddles}
The source proves the following two-term expansion for every $x>0$:
\begin{align}\label{eq:Ytwo}
 Y_0(x)&=\frac{\sqrt\pi}\pi\left((\sin x-\cos x)x^{-1/2}
           -\frac{\sin x+\cos x}{8}x^{-3/2}\right)+E_Y(x),\\
 |E_Y(x)|&\le\frac2\pi\frac9{128}\sqrt\pi\,x^{-5/2}.\nonumber
\end{align}
It follows from the proved complex Laplace-ray representation of the literal
Schl\"afli integral, a second-order amplitude remainder, and evaluation of the
first two Laplace moments. Thus no large-argument Bessel expansion is left as
a premise. Expressing the sine/cosine terms as exponentials retains two leading
coefficients
\[
 c_+=-(1+\ii)/2,\qquad c_-=-(1-\ii)/2
\]
and two correction coefficients $(1-\ii)/2,(1+\ii)/2$, together with the
$1/8$ factor. The complete correction and remainder series are proved
absolutely integrable and are bounded by $O_\delta(G\log T)$ in the physical
width range. Their terms are not silently discarded.
\src{Stronger/Atkinson/NeumannTwoTermExpansion.lean; ZetaNeumannSeries.lean; AtkinsonCorrectionSeries.lean; ZetaAtkinsonTwoTermSource.lean}{abs_dfiBesselY0_sub_neumannTwoTerm_le; exists_zetaAtkinsonBesselMinus_twoTerm_log_bound; exists_zetaSquareLocalMean_le_atkinson_twoTerm}

For a real signed frequency $b$, the leading continuous carrier is
\[
 \Phi(T,b;x)=T\log x-2\pi x+4\pi b\sqrt x.
\]
Write $A=T/(2\pi)$ and
\[
 r_b=\frac{b+\sqrt{b^2+4A}}2,\qquad x_b=r_b^2.
\]
For $T>0$, $r_b>0$ and $r_b^2-br_b=A$. Differentiation gives
\[
 \Phi_x=T/x-2\pi+2\pi b/\sqrt x
         =\frac{2\pi}{x}(A-x+b\sqrt x).
\]
The positive solution of $x-b\sqrt x=A$ is uniquely $x_b$, so it is the
unique positive saddle. Its second derivative is
$-\pi(2r_b-b)/r_b^3<0$. The exact paired identities are
\[
 r_{-b}=r_b-b,\quad r_br_{-b}=A,\quad x_bx_{-b}=A^2,\quad
 \log(x_b/A)=2\arsinh\!\left(\frac b{2\sqrt A}\right).
\]
All are obtained from the positive quadratic root and positive logarithm;
there is no choice of a second branch.
\src{Stronger/Atkinson/ZetaAtkinsonSaddle.lean; AtkinsonSaddleNormalization.lean}{zetaAtkinsonPhase_stationary_iff; zetaAtkinsonPhase_secondDeriv_saddle_neg; atkinsonSaddleRoot_mul_neg; zetaAtkinsonSaddle_log_argument}

The exact amplitude for power $\alpha$ is
\[
 a_\alpha(T,G,L;x)=x^{-\alpha}W(\log x-\ell_T)
       \vartheta_{T,G,L}(x)\chi_TQ(T,G;\log x-\log A),
\]
on positive support. Define its carrier integral by
\[
 \mathcal A_\alpha(T,G,L;b)
       =\int_0^\infty x^{-1/2}a_\alpha(T,G,L;x)e^{\ii\Phi(T,b;x)}\dd x.
\]
The extra $x^{-1/2}$ belongs to the original divisor test; it is precisely
canceled by the Jacobian under $x=y^2$.
Put $h_b=1+A/r_b^2>1$. The stationary term is exactly
\begin{equation}\label{eq:stationaryterm}
 \mathcal S_\alpha(T,G,L;b)=
 2a_\alpha(T,G,L;x_b)e^{\ii\Phi(T,b;x_b)}
             \frac{e^{-\ii\pi/4}}{\sqrt{2h_b}}.
\end{equation}
The $2$ is the root-coordinate Jacobian combined with the source's root
amplitude convention. The Fresnel value is proved from a damped Gaussian and
an Abel limit; its tails are estimated before removing the damping.
\src{Stronger/Atkinson/AtkinsonPowerWeight.lean; AtkinsonRootIntegral.lean; AtkinsonQuadraticKernel.lean; AtkinsonStationaryMain.lean; FresnelEvaluation.lean}{atkinsonPowerWeight; atkinsonPowerIntegral; atkinsonStationaryMain; fresnelGaussianValue_eq_phase}

\begin{lemma}[Uniform stationary error above the fourth-root width]\label{lem:stationaryerror}
For each fixed real $\alpha$ there is $C_\alpha>0$ such that, if
\[
 T\ge1,\quad T^{1/4}\le G\le\sqrt T,\quad L\ge1,\quad8L\le G,
 \quad |b|\le\sqrt T/100,
\]
then
\[
 |\mathcal A_\alpha(T,G,L;b)-\mathcal S_\alpha(T,G,L;b)|
 \le C_\alpha G\sqrt G\,T^{-\alpha-1/4}.
\]
Both $b=\sqrt n$ and $b=-\sqrt n$ satisfy the same estimate when
$10000n\le T$.
\end{lemma}
\begin{proof}
Set $x=y^2$. The root phase is
$2\pi(2A\log y-y^2+2by)$, with saddle $r_b$ and quadratic part
$\Phi(T,b;x_b)-2\pi h_b(y-r_b)^2$.
The source proves compact $C^2$ bounds for the root amplitude: its norm is at
most a constant times $GT^{-\alpha}$, and first and second derivatives have
scales $G/\sqrt T$ and $(G/\sqrt T)^2$. These follow by differentiating the
fixed Mellin profile, the smooth band cutoff, and the actual quadratic
Gaussian; bounds remain uniform on the support, rather than depending on
$n$ or the chosen saddle.

For a symmetric radius $h$ around $r_b$, the exact logarithmic Taylor
expansion has cubic correction $2T(y-r_b)^3/(3r_b^3)$ and a bounded
fourth-order remainder. Expand the amplitude to its linear term and the
exponential phase to its cubic term. The linear amplitude and the cubic phase
correction are odd around $r_b$, while the quadratic exponential is even;
their two integrals vanish exactly. The remaining local error is bounded,
with amplitude norm bound $M$ and derivative scale $R$, by
\[
 M\left(2R^2h^3+4RT h^5/r_b^3+8Th^5/r_b^4+8T^2h^7/r_b^6\right).
\]
This is obtained by bounding the three residual products before integration
and multiplying by the interval length $2h$. The nonstationary portions are
handled by the proved first-derivative/integration-by-parts estimates, and
the finite quadratic window differs from its Fresnel value by an explicit
reciprocal-radius tail.

For $|b|\le\sqrt T/100$, the source proves $r_b\asymp\sqrt T$ with explicit
window margins. Choose
$X=T^{1/4}/\sqrt G$ and actual local radius $h=X/12$.
The width assumptions give
$X\ge1$, $GX^2=\sqrt T$, $X^4=T/G^2\le\sqrt T$.
After inserting $R=G/\sqrt T$, the combined reciprocal-radius and symmetric
error is bounded by
\[
 \frac4{h\pi}+\frac{4G^2h^3}T+\frac{216Gh^5}T
       +\frac{1296h^5}T+\frac{11664h^7}T\le20/X.
\]
Each of the last four terms is bounded by $1/X$ from the displayed balances;
the first is at most $16/X$ because $\pi>3$.
Multiplying by $C_\alpha GT^{-\alpha}$ gives
$C_\alpha GT^{-\alpha}/X=C_\alpha G\sqrt G T^{-\alpha-1/4}$.
The small-frequency bound for both signs of $\sqrt n$ is exactly the
consequence of $10000n\le T$.
\end{proof}
\src{Stronger/Atkinson/StationaryOddRemainder.lean; AtkinsonSymmetricStationary.lean; AtkinsonSymmetricBalance.lean; AtkinsonFourthRootRadius.lean}{integral_atkinsonRootQuadratic_odd; IntervalC2Bound.atkinsonLocalSymmetric; exists_atkinsonPowerIntegral_symmetric_balance; exists_atkinsonPowerIntegral_fourthRoot_pair}

Define the complete stationary leading sum
\begin{equation}\label{eq:stationarysum}
 \begin{split}
 \mathcal S(T,G,L)&=\sum_{n\ge1}d(n)(-2\pi)
 \frac{\sqrt\pi}\pi(4\pi)^{-1/2}n^{-1/4}\\
 &\quad\times\left(c_+\mathcal S_{1/4}(T,G,L;\sqrt n)
      +c_-\mathcal S_{1/4}(T,G,L;-\sqrt n)\right).
 \end{split}
\end{equation}
For $8L\le G$, both saddle amplitudes vanish when
$n\ge N_*(T,G,L)=\lceil36T(L/G)^2\rceil$; hence this sum is a finite sum
with exactly the displayed terms. The source proves the cutoff condition from
the saddle logarithms and the support of $\vartheta$, before taking the sum.
The original leading carrier is also truncated at this cutoff with a proved
tail bound. Lemma~\ref{lem:stationaryerror} at $\alpha=1/4$ and the divisor
prefix bound $\sum_{n<N}d(n)n^{-1/4}\ll_\theta N^{3/4+\theta}$ give the
finite sum error
\[
 O_\theta\bigl(G\sqrt G\,T^{-1/2}N_*^{3/4+\theta}\bigr).
\]
When $G\ge T^{1/4+\kappa}$ and $L=\log T$, a smaller positive loss depending
on $\delta,\kappa$ makes this $O(G)$; the strict $\kappa$ margin absorbs the
logarithms. The source also proves that the complete two-term Neumann source
differs from this complete stationary sum by $O_{\delta,\kappa}(G)$.
Combining with \eqref{eq:localsource} yields
\begin{equation}\label{eq:stationarylocal}
 I(T,G)\le2e\Rea\mathcal S(T,G,\log T)+C_{\delta,\kappa}G\log T
\end{equation}
for sufficiently large $T$ and
$T^\delta\le G\le T^{1/2-\delta}$, $G\ge T^{1/4+\kappa}$.
\src{Stronger/Atkinson/AtkinsonStationarySeries.lean; AtkinsonFiniteStationaryError.lean; AtkinsonTwoTermStationary.lean; AtkinsonNativeLocalMean.lean}{atkinsonStationaryLeadingSum_eq_finite; exists_norm_atkinsonLeadingFiniteSum_sub_stationary_le; exists_tsum_zetaAtkinsonTwoTerm_sub_stationary_above_fourthRoot; localMeanStationaryInput_native}

\subsection{Exact signed normalization and prefix control}
The phase evaluated at the two saddles is expressed with
\[
 f(T,n)=2T\arsinh\sqrt{\frac{\pi n}{2T}}
       +\sqrt{2\pi nT+\pi^2n^2}-\frac\pi4,
 \qquad c(T)=T\log(T/(2\pi))-T.
\]
It satisfies
\[
 \Phi(T,\sqrt n;x_{\sqrt n})=c(T)+\pi n+f(T,n)+\pi/4,
\]
and the negative saddle gives $c(T)+\pi n-f(T,n)-\pi/4$.
After multiplying by the Fresnel factor $e^{-\ii\pi/4}$, the positive
carrier is $e^{\ii f}$ while the negative carrier is
$-\ii e^{-\ii f}$. Neither is replaced by an unsigned absolute main term.

To display the full signed weights define
\begin{gather*}
 C_*(T,b)=\frac1{\sqrt2}(b^2+2T/\pi)^{-1/4},\qquad
 c_n(T)=n^{-1/4}C_*(T,\sqrt n),\\
 U(T,G;n)=Q\!\left(T,G;2\arsinh\sqrt{\pi n/(2T)}\right),\\
 R_\pm(T,G,L;n)=W(\log x_{\pm\sqrt n}-\ell_T)
                         \vartheta_{T,G,L}(x_{\pm\sqrt n}),\\
 w_\pm(T,G,L;n)=c_n(T)U(T,G;n)R_\pm(T,G,L;n),\\
 E_\pm(T,n)=d(n)(-1)^n e^{\pm\ii f(T,n)}.
\end{gather*}
Then the exact identity for \eqref{eq:stationarysum} is
\begin{equation}\label{eq:signedstationary}
 \mathcal S(T,G,L)=(1+\ii)\chi_Te^{\ii c(T)}
       \left(\sum_{n<N_*}w_+(T,G,L;n)E_+(T,n)
                    -\sum_{n<N_*}w_-(T,G,L;n)E_-(T,n)\right).
\end{equation}
It follows by algebraic substitution of the evaluated saddle formulas and the
fourth-root normalization
$2\sqrt\pi(4\pi)^{-1/2}n^{-1/4}=n^{-1/4}$.
The two residual weights are different functions at different saddles;
$E_- =\overline{E_+}$ does not equate the weights. The common phase has modulus
at most two.
\src{Stronger/Atkinson/AtkinsonEvaluatedPhases.lean; AtkinsonCoefficientVariation.lean; AtkinsonSaddleProfileNormalization.lean; AtkinsonSignedPhaseSeries.lean}{atkinsonStationaryMain_neg_sqrt; atkinsonFourthRootCoefficient_eq; atkinsonStationaryLeadingSum_eq_signed}

For a finite sequence $w_i$, the source's
\lean{FiniteVariationBound} consists of a uniform bound $|w_i|\le M$ for
$0\le i\le N$ and a separate bound
$\sum_{i<N}|w_{i+1}-w_i|\le M$, with $M\ge0$.
The finite summation-by-parts identity bounds
$|\sum_{i<N}w_iE_i|$ by $2M$ times the maximum of the \emph{actual}
prefix sums $|\sum_{i<j}E_i|$, $0\le j\le N$.
For the two weights in \eqref{eq:signedstationary}, the source proves the
variation bounds
\[
 \operatorname{Var}_{m,N}(w_\pm)
 \le C G T^{-1/4}m^{-1/4}e^{-G^2m/(12T)}
\]
when $m\ge1$ and $10000(m+N)\le T$. The real fourth-root coefficient is
positive antitone, so its variation telescopes. The Gaussian increment bound
is derived by differentiating $Q$, and the two saddle residuals have uniform
variation by differentiating the actual compact profiles. Multiplying the
three variation bounds yields the displayed estimate.

Let
\[
 P_{m,N}(T)=\max_{0\le j\le N}
           \left|\sum_{i=0}^{j-1}d(m+i)(-1)^{m+i}e^{\ii f(T,m+i)}\right|,
\]
and define the undamped dyadic prefix bound
\[
 \mathscr B(T,N)=\sum_{0\le j<\lceil\log_2N\rceil}
                     (2^j)^{-1/4}P_{2^j,2^j}(T).
\]
For $N=0,1$ the source uses the natural logarithm-cutoff convention; the empty
sum is zero, and the zero-indexed stationary term is zero. Truncated dyadic
pieces partition every finite prefix exactly. Summation by parts on each,
then enlargement of its prefix maximum to the full block, gives
\begin{equation}\label{eq:prefixsource}
 |\mathcal S(T,G,\log T)|\le C_\delta GT^{-1/4}
           \mathscr B(T,N_*(T,G,\log T)).
\end{equation}
The damped factors are first kept; only their inequality $e^{-G^2m/(12T)}\le1$
is used when passing to this undamped bound. Equations
\eqref{eq:stationarylocal} and~\eqref{eq:prefixsource} provide the actual
local-source-to-prefix interface.
\src{Stronger/Atkinson/FiniteWeightVariation.lean; AtkinsonMainVariation.lean; AtkinsonStationaryBlockBound.lean; AtkinsonDampedDyadic.lean; AtkinsonStationaryDyadic.lean; AtkinsonStationarySourceAssembly.lean}{exists_finiteVariationBound_atkinsonMainWeights; exists_norm_atkinsonStationarySum_le_undamped; localMeanSourceToPrefix_of_stationary}

\section{From signed prefixes to physical local-mean packets}
\label{sec:packets}
\subsection{The actual maximal-prefix Gram construction}
For a positive dyadic block $[m,m+N)$, the coefficients are
$\alpha_i=d(m+i)(-1)^{m+i}$ and its unit phase vector at height $t$ is
$v_t(i)=e^{\ii f(t,m+i)}$. At each height select a maximizing prefix length
$j_t\in\{0,\ldots,N\}$ for $P_{m,N}(t)$. Define the masked vector
$v_t^{j_t}(i)=\mathbf1_{i<j_t}v_t(i)$.
The exact Gram entry of two such vectors is
\[
 \sum_{i<N}\overline{v_t^{j_t}(i)}v_u^{j_u}(i)
   =\sum_{i<\min(j_t,j_u)}e^{\ii(f(u,m+i)-f(t,m+i))}.
\]
Thus each pair keeps the minimum of its two own maximizing prefixes. There is
no single maximizing prefix imposed on all heights. If
\[
 G_{m,N}(t,u)=\max_{0\le j\le N}
       \left|\sum_{i<j}e^{\ii(f(u,m+i)-f(t,m+i))}\right|,
 \quad E_{m,N}=\sum_{i<N}d(m+i)^2,
\]
finite Hilbert-space Cauchy--Schwarz applied to the sum of phase-aligned masked
vectors yields
\begin{equation}\label{eq:Gram}
 \left(\sum_{t\in W}P_{m,N}(t)\right)^2
          \le E_{m,N}\sum_{t,u\in W}G_{m,N}(t,u).
\end{equation}
Indeed choose a unit scalar at each height to make its coefficient-vector inner
product real and nonnegative. Cauchy--Schwarz bounds the square of their sum by
the coefficient norm squared times the summed-vector norm squared; expanding
that norm and taking absolute values gives \eqref{eq:Gram}. Zero prefixes or
zero values cause no difficulty. The source proves this finite lemma before
using any analytic cancellation estimate.
\src{Stronger/Atkinson/AtkinsonPrefixVectors.lean; AtkinsonPrefixMax.lean; FiniteWeightedGram.lean; AtkinsonDyadicGramBudget.lean}{atkinsonMaskedPhaseVector_gram; sum_atkinsonPhaseBlockMax_sq_le_gramMax; sum_atkinsonUndampedDyadicPhaseBound_sq_le_budget}

Finite Cauchy--Schwarz over dyadic blocks gives, with $J=\lceil\log_2N\rceil$,
\begin{equation}\label{eq:dyadicGram}
 \left(\sum_{t\in W}\mathscr B(t,N)\right)^2
 \le J\sum_{j<J}(2^j)^{-1/2}E_{2^j,2^j}
                        \sum_{t,u\in W}G_{2^j,2^j}(t,u).
\end{equation}
The right side is exactly the divisor-weighted dyadic Gram budget of the source.
The divisor subpower estimate bounds
$(2^j)^{-1/2}E_{2^j,2^j}\ll_\eta(2^j)^{1/2+\eta}$ for every $\eta>0$.

\subsection{First- and second-derivative bounds for each prefix}
Treat the summation index as a positive real variable $x$. The exact phase
slope and curvature are
\[
 f_x(T,x)=\sqrt{2\pi T/x+\pi^2},\qquad
 f_{xx}(T,x)=-\frac{\pi T}{x^2\sqrt{2\pi T/x+\pi^2}}.
\]
For $0<x\le u\le t\le2u$, the positive difference
$f(u,x)-f(t,x)$ has positive second derivative. Differentiating curvature
with respect to height and integrating from $u$ to $t$ gives explicit upper
and lower scales. On the entire block $m\le x\le m+N+2$ they are
\begin{gather*}
 \lambda=\frac{\pi(t-u)}{2(m+N+2)^2
        \sqrt{(2\pi+\pi^2)2u/m}},\\
 \lambda'=\frac{\pi(t-u)}{m^2\sqrt{2\pi u/(m+N+2)}}.
\end{gather*}
Both are positive for $t>u$, $m\ge1$. A second-order mean-value theorem
puts every unit-spaced second difference between $\lambda$ and $\lambda'$.
The finite van der Corput B-process, proved in the imported local
\lean{GuthMaynard.SecondDerivative} development, therefore bounds each
prefix by
\[
 \left(\frac{N\lambda'}{2\pi}+2\right)
 \left(\frac{2\pi}{\sqrt\lambda}
       +2\left(\frac{\sqrt\lambda}{\lambda}+1\right)\right).
\]
A shorter prefix uses the same block curvature parameters; its length factor
is merely enlarged to $N$. The zero prefix is handled separately.

For near height gaps, the first difference of $f(t,x)-f(u,x)$ is monotone
and positive. The exact first-derivative scales place it in
$[\delta_*,\pi]$ when $|t-u|\le\sqrt{Hm}$ and $t,u\in[H,2H]$,
where $\delta_*>0$ is the proved lower difference bound. The Kusmin--Landau
one-period estimate gives $2\pi/\delta_*$. The source takes the minimum of
this bound, the B-process bound, and the triangle bound $N$. For the dyadic
case $N=m$, its orientation-free majorant satisfies
\begin{align}\label{eq:nearfar}
 G_{m,m}(t,t)&=m,\\
 G_{m,m}(t,u)&\le60\frac{\sqrt{Hm}}{|t-u|}
                  &&(0<|t-u|\le\sqrt{Hm}),\nonumber\\
 G_{m,m}(t,u)&\le2000\sqrt{\frac{m|t-u|}{\sqrt{Hm}}}
                  &&(|t-u|\ge\sqrt{Hm}).\nonumber
\end{align}
These statements require the block-plus-two endpoint at most the smaller
height and the larger height at most twice the smaller. The native physical
cutoff guarantees both conditions simultaneously. In the far range the source
uses the triangle branch when curvature is large, so there is no extra linear
gap term left in \eqref{eq:nearfar}. More explicitly, put
$d=|t-u|$, $R=\sqrt{Hm}$ and $Q=mR$.
The actual curvature bounds give $\lambda\ge d/(320Q)$,
$\lambda'\le4d/Q$. If $R\le d\le Q$, the first B-process factor is
at most $6d/R$, while the second is at most $202\sqrt{Q/d}$.
Their product is at most $1212\sqrt{md/R}\le2000\sqrt{md/R}$.
If $d>Q$, the triangle bound $m\le\sqrt{md/R}$ is stronger.
These two cases prove the far estimate without an upper restriction on the gap.
\src{Stronger/Atkinson/AtkinsonIndexCurvature.lean; AtkinsonIndexScales.lean; AtkinsonIndexBProcess.lean; AtkinsonPrefixCancellation.lean; AtkinsonNearGap.lean; AtkinsonFarGap.lean}{hasDerivAt_deriv_atkinsonIndexRealPhase; atkinsonIndexPositiveDifference_B_process; atkinsonPrefixGramMax_le_firstDerivative; atkinsonPrefixGapMajorant_le_near; atkinsonPrefixGapMajorant_le_far}

\begin{lemma}[Physical prefix packet budget]\label{lem:prefixpacket}
For $\delta,\eta>0$ there are $C>0$ and $H_0\ge40000$ such that, for
$H\ge H_0$, $H^\delta\le G$, $L\ge0$, a $G$-separated finite
$W\subset[H,2H]\cap[A,A+L]$, and
$N=N_*(2H,G,\log(2H))$, $R=|W|$,
\[
 \left(\sum_{t\in W}\mathscr B(t,N)\right)^2\le C\mathcal P_\eta(H,G,L,N,R),
\]
where $J_N=\lceil\log_2N\rceil$ and
\begin{align*}
 \mathcal P_\eta(H,G,L,N,R)
 &=J_N^2\left[R N^{3/2+\eta}
  +\frac{120R\sqrt H}{G}N^{1+\eta}
       \mathsf H_{\lceil\sqrt{HN}/G\rceil}\right.\\
 &\hspace{42mm}\left.+2000R^2H^{-1/4}\sqrt L\,N^{3/4+\eta}\right],\\
 \mathsf H_k&=\sum_{j=1}^k1/j.
\end{align*}
\end{lemma}
\begin{proof}
Eventually $1200\log(2H)\le H^{\delta/2}\le G$ after the harmless doubled
height adjustment. The explicit ceiling cutoff then satisfies
$2N+2\le H$ and $N\le H$. Hence every dyadic block in
\eqref{eq:dyadicGram} satisfies the prefix geometry needed for
\eqref{eq:nearfar}. The diagonal contributes $Rm$. For fixed $t$, order the
$G$-separated points on either side of $t$. The $j$th has distance at least
$jG$, so the near reciprocal sum is at most
$120\sqrt{Hm}\,\mathsf H_{\lceil\sqrt{Hm}/G\rceil}/G$.
Since every pair gap is at most $L$, the far double sum is at most
$2000R^2\sqrt{mL/\sqrt{Hm}}$.
Multiply these three contributions by the coefficient factor
$Cm^{1/2+\eta}$. They become respectively
$CRm^{3/2+\eta}$,
$C(120R\sqrt H/G)m^{1+\eta}\mathsf H_{\lceil\sqrt{Hm}/G\rceil}$,
and $C2000R^2H^{-1/4}\sqrt Lm^{3/4+\eta}$.
Each is nondecreasing with $m$ on positive indices; use $m\le N$.
The outer dyadic Cauchy--Schwarz factor and the number of blocks give
$\lceil\log_2N\rceil^2$. This proves the stated budget, including the
empty-set case, for which all terms are zero.
\end{proof}
\src{Stronger/Atkinson/AtkinsonNearGapBudget.lean; AtkinsonPowerGapBudget.lean; AtkinsonPhysicalCutoff.lean; AtkinsonLocalizedGram.lean}{exists_atkinsonPhysicalCutoff_prefix_geometry; exists_sum_atkinsonUndampedDyadicPhaseBound_sq_le_power}

\begin{proposition}[Native actual local-mean packet]\label{prop:localpacket}
For every $\delta,\kappa,\nu>0$ there exist positive $C,D$ and $H_0\ge40000$
such that the following holds. Let $H\ge H_0$, $G>0$, $L\ge0$ and a finite
$G$-separated set $W$ satisfy, for every $t\in W$,
\[
 H\le t\le2H,\quad t^\delta\le G\le t^{1/2-\delta},\quad
 t^{1/4+\kappa}\le G,\quad A\le t\le A+L.
\]
Set $E_G(t)=\max(0,I(t,G)-CG\log t)$. Then
\begin{equation}\label{eq:localpacket}
 \left(\sum_{t\in W}E_G(t)\right)^2
       \le DH^\nu\left(\frac{|W|H}{G}+|W|^2\sqrt{GL}\right).
\end{equation}
\end{proposition}
\begin{proof}
Equations \eqref{eq:stationarylocal} and~\eqref{eq:prefixsource} give
$E_G(t)\le CGt^{-1/4}\mathscr B(t,N_*(t,G,\log t))$.
The cutoff is monotone in $t\ge1$ and the undamped prefix bound monotone in its
cutoff. Thus $t\in[H,2H]$ yields
\[
 E_G(t)\le CGH^{-1/4}\mathscr B(t,N),
 \qquad N=N_*(2H,G,\log(2H)).
\]
Apply Lemma~\ref{lem:prefixpacket} with loss $\nu/2$ and square the factor
$CGH^{-1/4}$.

To evaluate the physical budget, the ceiling bound gives
$N\le74H(\log(2H))^2/G^2$, since $G\le\sqrt{2H}$.
The diagonal index power then satisfies
\[
 G^2H^{-1/2}N^{3/2}\le5476(H/G)(\log(2H))^3.
\]
The near power becomes
$(\sqrt H/G)(G^2H^{-1/2}N)=74(H/G)(\log(2H))^2$.
The harmonic factor is at most $2\log(2H)$ and the dyadic-count factor at
most $((1/\log2+1)\log(2H))^2$.
The far power satisfies
\[
 H^{-1/4}\sqrt L\,G^2H^{-1/2}N^{3/4}
       \le74\sqrt{GL}(\log(2H))^3.
\]
The factor $N^{\nu/2}\le H^{\nu/2}$ and the remaining fifth power of the
logarithm is eventually at most $H^{\nu/2}$. Substitution into all three
terms proves \eqref{eq:localpacket}. The constants are chosen before
$H,G,A,L,W$ and the source interface has been discharged by the native
stationary producer. A nonempty packet gives $G\le\sqrt{2H}$ from one of its
points; an empty packet has zero sum and is dealt with directly.
\end{proof}
\src{Stronger/Atkinson/AtkinsonPhysicalBudget.lean; AtkinsonSourceAssembly.lean; AtkinsonStationarySourceAssembly.lean; AtkinsonNativeLocalMean.lean}{exists_atkinsonPhysicalGapBudget_le_twoTerm; localMeanPacketInput_of_sourceToPrefix; localMeanPacketInput_of_stationary; localMeanStationaryInput_native}

\section{High point values, finite occupancy and the twelfth moment}
\label{sec:peaks}
\subsection{Absorption in a short physical interval}
\begin{lemma}[Finite local-to-global absorption]\label{lem:absorption}
Let $A,H,G,Y>0$, and a finite $W\subset[H,2H]$ have nonnegative values
$f(t)\ge Y$. Suppose every $U\subset W$ localized in an interval of length
\[
 L_*=(Y^2/(4A))^2/G
\]
satisfies
$\bigl(\sum_{t\in U}f(t)\bigr)^2\le A(|U|H/G+|U|^2\sqrt{GL_*})$.
Then
\[
 |W|\le\frac{2AH}{GY^2}+\frac{32A^3H^2}{Y^6}.
\]
\end{lemma}
\begin{proof}
$\sqrt{GL_*}=Y^2/(4A)$, so $2A\sqrt{GL_*}\le Y^2$.
For $R=|U|>0$, use $\sum f\ge RY$ in the local packet and absorb its
quadratic term to get $R\le2AH/(GY^2)$. The empty fiber also satisfies this.
Partition $[H,2H]$ by the natural floor bins of $(t-H)/L_*$;
at most $\lfloor H/L_*\rfloor+1$ fibers are occupied. This convention also
captures a point on the closed right boundary. Sum the fiber cardinality
bounds and use $\lfloor H/L_*\rfloor\le H/L_*$. Expanding gives exactly
$2AH/(GY^2)+32A^3H^2/Y^6$.
\end{proof}
\src{Stronger/AtkinsonCardinalityAbsorption.lean; AtkinsonHeightCover.lean; AtkinsonGlobalCardinality.lean}{atkinsonAbsorptionLength; atkinson_card_le_of_packet; atkinson_card_le_of_local_packets}

Applying this lemma to Proposition~\ref{prop:localpacket}, with its loss
$\nu/3$, yields the actual local-mean-excess count
\begin{equation}\label{eq:excesscount}
 |W|\le D H^\nu\left(\frac H{GY^2}+\frac{H^2}{Y^6}\right)
 \quad(E_G(t)\ge Y).
\end{equation}
The cube of $H^{\nu/3}$ is $H^\nu$; the first term is enlarged to that same
loss. Thus no unseen exponent is spent in the global height cover.

\subsection{Clusters of one-separated peaks}
Let $W\subset[H,2H]$ be one-separated and satisfy $\zabs(t)\ge V>0$.
For a width $G$, partition by
$n=\lfloor(t-H)/(G/2)\rfloor$. The center assigned to bin $n$ is its
\emph{left} endpoint $c_n=H+nG/2$, not its midpoint. Its cluster lies in
$[c_n,c_n+G/2]$, hence also in $[c_n-G/2,c_n+G/2]$.
The nonempty cluster centers remain in $[H,2H]$. Centers with equal parity
are $G$-separated; split into even and odd bins before using
\eqref{eq:excesscount}.

Set $\ell=\log(3H)$, and let $P>0$ be the point-mean constant in
Lemma~\ref{lem:equation44}. If $2\ell^2\le G\le H$, that lemma gives, for
cluster occupancy $r_n$,
\[
 V^2r_n\le P\ell(r_n+4I(c_n,G)).
\]
If $V^2\ge2P\ell$, then
\[
 I(c_n,G)\ge2r_n A_V,\qquad A_V=V^2/(16P\ell).
\]
If the local-mean error $CG\ell\le A_V$, then for every integer $m\ge1$
with $r_n\ge m$, the excess at $c_n$ is at least $mA_V$.
Apply \eqref{eq:excesscount} to the two parity classes of centers. The result
is
\[
 \#\{n:r_n\ge m\}\le2D H^\nu
 \left(\frac H{G(mA_V)^2}+\frac{H^2}{(mA_V)^6}\right)
 \le\frac{2D H^\nu}{m^2}
       \left(\frac H{GA_V^2}+\frac{H^2}{A_V^6}\right).
\]
The exact integer layer-cake identity $\sum_nr_n=\sum_{m\ge1}\#\{n:r_n\ge m\}$
and $\sum_{m\ge1}m^{-2}\le2$ therefore give
\begin{equation}\label{eq:widthcount}
 |W|\le D' H^\nu\left(\frac H{GA_V^2}+\frac{H^2}{A_V^6}\right).
\end{equation}
This occupancy step restores \emph{all} peaks, not merely one representative
per cluster.
\src{Stronger/Peaks/Clusters.lean; ClusterCount.lean; FiniteOccupancy.lean}{pointClusterCenter; pointClusterCenters_separated; exists_pointCluster_superlevel_entry; exists_pointCluster_superlevel_count; sum_nat_occupancy_le_of_superlevels; exists_pointValue_card_le_with_width}

Choose $K=16PC+1$ and
\[
 G=\frac{V^2}{K\ell^2}.
\]
The error absorption is automatic: $CG\ell\le V^2/(16P\ell)=A_V$.
When
\[
 K\ell^2(2H)^{1/4+\kappa}\le V^2
        \le K\ell^2H^{1/2-\delta},\qquad0<\delta\le1/4,
\]
all physical width conditions hold simultaneously throughout $[H,2H]$;
the lower bound also makes $2\ell^2\le G$ eventually.
Substitution into \eqref{eq:widthcount} gives
\begin{equation}\label{eq:sourcepeakcount}
 |W|\le D H^\nu\left(\frac{H\ell^4}{V^6}
                           +\frac{H^2\ell^6}{V^{12}}\right).
\end{equation}
The factors of $K$ and $16P$ are explicitly absorbed into $D$; the source
checks the identity before using it.
\src{Stronger/Peaks/WidthRanges.lean}{pointValueWidth_error_absorption; pointValueWidth_count_identity; exists_pointValue_card_le_source_range}

\begin{proposition}[Native high-value finite-set bound]\label{prop:peakcard}
For every $\varepsilon>0$, eventually in $H$, for all $V>0$ with
$V\ge H^{1/8+\varepsilon}$ and every one-separated finite
$W\subset\{t\in[H,2H]:\zabs(t)\ge V\}$,
\[
 |W|V^{12}\le H^{2+\varepsilon}.
\]
\end{proposition}
\begin{proof}
Choose $\eta=\min(\varepsilon/16,1/1000)$, use
$\delta=1/48$, $\kappa=\eta$, $\nu=\eta$ in
\eqref{eq:sourcepeakcount}. The lower source range follows eventually from
$V^2\ge H^{1/4+2\varepsilon}\ge H^{1/4+2\eta}$, because the extra
positive $\eta$ power absorbs $K\ell^2 2^{1/4+\eta}$.
If $W$ is empty there is nothing to prove. Otherwise take a point in $W$.
Proposition~\ref{prop:weyl} gives $V\le H^{1/6+\eta}$. Consequently
$V^2\le H^{1/3+2\eta}\le K\ell^2 H^{1/2-1/48}$ eventually, so the
upper source range also holds uniformly.
Multiply \eqref{eq:sourcepeakcount} by $V^{12}$:
\[
 |W|V^{12}\le D H^\eta(H\ell^4V^6+H^2\ell^6)
       \le D\ell^4H^{2+7\eta}+D\ell^6H^{2+\eta}.
\]
Each logarithm and constant is absorbed into the spare positive loss, giving
$H^{2+8\eta}\le H^{2+\varepsilon}$.
The choices and thresholds depend on $\varepsilon$ and the two native
analytic producers, but precede $H,V,W$. Substituting the proved point-mean
and stationary packet producers removes both explicit input hypotheses.
\end{proof}
\src{Stronger/Peaks/WidthRanges.lean; NativePeaks.lean}{exists_pointValue_twelfth_weighted_card_le_of_inputs; pointValue_peak_card_of_inputs; pointValue_peak_card_native}

\begin{lemma}[From every finite separated set to full measure]\label{lem:peakmeasure}
Put $S(H,V)=\{t:H\le t\le2H,\ \zabs(t)\ge V\}$.
Under the eventual conditions of Proposition~\ref{prop:peakcard},
\[
 |S(H,V)|_{\rm Leb}\le2H^{2+\varepsilon}/V^{12}.
\]
\end{lemma}
\begin{proof}
A uniform cardinality bound $|W|\le B$ on every one-separated finite subset
of a set $S$ implies that the possible cardinalities have a largest value:
they lie in the finite set $\{0,\ldots,\lceil B\rceil\}$.
Choose a one-separated subset of that maximal size. If a point of $S$ were
farther than one from every chosen point, it could be inserted to increase the
size. Thus closed radius-one balls about the chosen points cover $S$.
Their total length is at most $2B$. Take $B=H^{2+\varepsilon}/V^{12}$.
Continuity of $\zabs$ makes $S(H,V)$ closed and measurable. This proves a
measure bound for the full superlevel set, not just for a sampled subset.
\end{proof}
\src{Stronger/SeparatedVolume.lean; NativePeaks.lean}{exists_unitSeparated_closedBall_cover_of_card_bound; volume_le_two_mul_of_separated_card_bound; pointValue_volume_native}

\subsection{Layer cake, the high/low split and the physical interval}
\begin{lemma}[Truncated high moment]\label{lem:highmoment}
Suppose $V>0$, $V^{12}\le M$, $C\ge0$, $\zabs(t)^{12}\le M$ on $S(H,V)$,
and $|S(H,U)|_{\rm Leb}\le C/U^{12}$ for every $U\ge V$.
Then
\[
 \int_{S(H,V)}\zabs(t)^{12}\dd t
       \le C\bigl(1+\log(M/V^{12})\bigr).
\]
\end{lemma}
\begin{proof}
For the restriction to $S(H,V)$ the distribution of $f(t)=\zabs(t)^{12}$
is bounded by $C/\max(V^{12},s)$ at level $s>0$. If $s\le V^{12}$ use
the whole-set bound; if $s>V^{12}$ use the level
$U=s^{1/12}\ge V$. Positivity and monotonicity of the twelfth power justify
this change of level. The function is integrable on this bounded measurable
set by continuity of $\zabs$. Apply layer cake up to $M$ and integrate:
\[
 \int_0^M\frac C{\max(V^{12},s)}\dd s
       =C+C\int_{V^{12}}^M\frac{\dd s}s
       =C\bigl(1+\log(M/V^{12})\bigr).
\]
\end{proof}
\src{Stronger/TruncatedLayerCake.lean; CriticalMoment.lean}{integral_le_log_of_truncated_tail; pointValue_power_tail_measure_le; zeta_twelfth_high_integral_le_log}

\begin{proof}[Proof of Theorem~\ref{thm:moment}]
First prove the dyadic bound. Given $\varepsilon>0$ put $q=\varepsilon/20$.
Proposition~\ref{prop:fourth} gives
$\int_H^{2H}\zabs^4\le C H^{1+q}$ eventually.
For the high-moment step use a smaller loss
$\eta=\min(q/2,1/100)$ in Lemma~\ref{lem:peakmeasure}, start at
$V_0=H^{1/8+\eta}$ and use the coarse upper truncation
$M=H^3$ from Proposition~\ref{prop:weyl}. Then
$1\le V_0^{12}\le H^3$ for $H\ge1$, and Lemma~\ref{lem:highmoment} gives
\[
 \int_{S(H,V_0)}\zabs(t)^{12}\dd t
 \le2H^{2+\eta}\bigl(1+\log(H^3/V_0^{12})\bigr)
 \le2H^{2+\eta}(1+3\log H)\le H^{2+q}
\]
eventually, since $\eta<q$. The set
$S(H,H^{1/8+q})$ is a subset of this starting high set, so its integral
obeys the same bound.

On the complement of $S(H,H^{1/8+q})$ within $[H,2H]$,
$\zabs^{12}\le H^{8(1/8+q)}\zabs^4$. Thus
\begin{align*}
 \int_H^{2H}\zabs(t)^{12}\dd t
 &\le H^{2+q}+H^{1+8q}\int_H^{2H}\zabs(t)^4\dd t\\
 &\le H^{2+q}+C H^{2+9q}
 \le(1+C)H^{2+\varepsilon}.
\end{align*}
This is the native dyadic producer, with a positive constant $1+C$.
All integrals over the high set and its complement are finite because they
are restrictions of a continuous function on a compact interval.

To pass to $[0,3T]$, let $B\ge1$ exceed the dyadic threshold and set
$r=2+\varepsilon\ge1$. For $f=\zabs^{12}$ choose
$A=|\int_0^Bf|+(1+C)+1>0$, so the dyadic constant $1+C\le A$.
Continuity proves that the compact initial integral is finite, and
$\int_0^Bf\le AB^r$. If
$\int_0^{2^mB}f\le A(2^mB)^r$, split at $2^mB$ and apply the dyadic
bound to the next interval. The sum is at most
$2A(2^mB)^r\le A(2^{m+1}B)^r$, because $2\le2^r$.
This induction is exactly the source's zero-origin adapter.
For $X\ge B$, choose $X\le2^mB\le2X$. Nonnegativity permits enlargement
to this dyadic endpoint, giving $\int_0^Xf\le A2^rX^r$ with no further loss.
Put $X=3T$ and use $(3T)^r=3^rT^r$; choose $K=A2^r3^r>0$.
The identity between interval and closed-set integrals gives the literal
$\int_{[0,3T]}$ endpoint. Finally continuity of $s(t)$ and analyticity of
$\zeta$ away from $1$ show $\zabs$ continuous on $\R$, hence all its natural
powers are integrable on every finite $[a,b]$. This proves both assertions.
\end{proof}
\src{Stronger/ConditionalMoment.lean; DyadicMoment.lean; NativeMoment.lean; Solution.lean}{zeta_twelfth_high_eventually_of_tail; zeta_twelfth_dyadic_of_fourth_and_tail; integral_zero_le_of_dyadic; zeta_twelfth_physical_native; actual_moment_integrable}

\section{The exact smooth zero detector}
\label{sec:detector}
We now reconstruct the zero-density consumer, using the analytic producers
already proved. Fix $3/4<\sigma<1$, and write
\begin{equation}\label{eq:parameters}
 a=\frac1{2\sigma},\qquad c=\frac3{4\sigma}=\frac32a,\qquad
 p=\frac{3(1-\sigma)}{2\sigma},\qquad X=T^\delta,\quad Y=T^c.
\end{equation}
Then $1/2<a<2/3$, $0<c<1$, $p>0$. Choose
$0<\delta\le1/8$ with $\delta<(1-c)/12$ before $T$ and any zero.
The source supplies, for example,
$\delta=\min(1/16,(1-c)/24)$.
The exact identities used later are
\begin{equation}\label{eq:exponentidentities}
 2(1-\sigma)c=p,\qquad\tfrac12+(3-4\sigma)a=p,\qquad
 2+(3-6\sigma)c=p-(1-c).
\end{equation}
They follow by substituting \eqref{eq:parameters} and clearing the positive
$\sigma$ denominator. They are balances of the two short terms and the long
saving, not asymptotic approximations.
\src{Stronger/Parameters.lean}{lowerExponent; smoothingExponent; densityExponent; exists_detector_exponent; short_exponent_identities; long_exponent_identity}

\begin{definition}[M\"obius mollifier and detector coefficients]
For a real $X$ define
\[
 M_X(s)=\sum_{1\le d\le\lfloor X\rfloor_+}\mu(d)d^{-s},\qquad
 b_X(n)=\sum_{d\mid n\atop d\le\lfloor X\rfloor_+}\mu(d),\qquad b_X(0)=0.
\]
Here $\lfloor X\rfloor_+$ is the natural floor used by Lean; for $X\ge0$
it is the ordinary natural floor. Put
\[
 \mathcal D(\rho;X,Y)=\sum_{n\ge1}b_X(n)n^{-\rho}e^{-n/Y},\qquad
 \mathcal T(\rho;X,Y)=\sum_{n\ge2}b_X(n)n^{-\rho}e^{-n/Y}.
\]
If $X\ge1$, then $b_X(1)=1$, $b_X(n)=0$ for $2\le n\le X$,
and $|b_X(n)|\le d(n)$. These follow from the exact M\"obius divisor
cancellation and $|\mu(d)|\le1$. Also
$|M_X(1/2+\ii t)|\le2\sqrt X$ for $X\ge0$, by the finite
$n^{-1/2}$ sum comparison.
\end{definition}
\src{MollifierCoefficients.lean; MollifierBound.lean; Stronger/InitialDetector.lean}{zetaMollifier; mollifierCoefficient_vanishes; mollifierCoefficient_norm_le; norm_zetaMollifier_le_two_sqrt; detectorTail}

\begin{lemma}[Exact exponentially smoothed contour identity]\label{lem:detectoridentity}
Let $\rho=\beta+\ii\gamma$ be a zero with $3/4\le\beta<1$, and $Y\ge1$.
The detector series is absolutely convergent and
\begin{align}\label{eq:detectoridentity}
 \mathcal D(\rho;X,Y)
 &=M_X(1)Y^{1-\rho}\Gamma(1-\rho)\\
 &\quad+\frac1{2\pi}\int_{\R}\Gamma(1/2-\beta+\ii u)
 \zeta(1/2+\ii(\gamma+u))M_X(1/2+\ii(\gamma+u))
 Y^{1/2-\beta+\ii u}\dd u.\nonumber
\end{align}
\end{lemma}
\begin{proof}
Mellin inversion of $e^{-n/Y}$ on the initial right line
$\Rea w=1/2$ gives the integral of
$\Gamma(w)\zeta(\rho+w)M_X(\rho+w)Y^w$.
This line has $\Rea(\rho+w)>1$, where the Dirichlet product equals the
series with coefficients $b_X(n)$. Absolute convergence and Gamma decay
justify exchanging the series and integral. Shift to
$\Rea w=1/2-\beta$. The apparent Gamma pole at zero is canceled by
$\zeta(\rho)=0$; its cleared product is analytic there. The remaining pole
is $w=1-\rho$, from zeta, with residue
$M_X(1)Y^{1-\rho}\Gamma(1-\rho)$. No other Gamma pole lies between these
lines. The finite rectangle identities prove this residue statement and the
integrability of both lines and vanishing of the horizontal edges. Taking the
limit gives \eqref{eq:detectoridentity}, since
$\rho+1/2-\beta+\ii u=1/2+\ii(\gamma+u)$.
The exponential series is absolutely convergent, using the coefficient bound
and the positive exponential damping; the source proves this before calling
it a detector.
\end{proof}
\src{DetectorMellin.lean; DetectorContour.lean; DetectorEstimates.lean; DetectorIdentity.lean}{smoothedDetector_eq_rightContour; detector_finite_rectangle_kernel; summable_zeta_detector_series; zeta_detector_identity}

\begin{lemma}[Uniform large tail]\label{lem:largetail}
For the choices \eqref{eq:parameters}, $0\le\delta\le1/8$, eventually in
$T$, simultaneously at every zero with
$\sigma\le\beta<1$ and $T\le|\gamma|\le2T$,
\[
 |\mathcal T(\rho;X,Y)|\ge1/2.
\]
\end{lemma}
\begin{proof}
Use Proposition~\ref{prop:weyl} with any fixed $0<\theta<1/48$ in the
integral term of \eqref{eq:detectoridentity}. The mollifier bound is
$2\sqrt X$. The shifted-height power is bounded by
$3T^{1/6+\theta}(1+|u|)$, and the Gamma strip has a uniform integrable
$(1+|u|)|\Gamma(1/2-\beta+\ii u)|$ majorant. Hence the integral term is
at most a fixed multiple of
\[
 \sqrt X\,Y^{1/2-\beta}T^{1/6+\theta}
       =T^{\delta/2+c(1/2-\beta)+1/6+\theta}
       \le T^{-1/48+\theta}.
\]
For the last inequality, $c(\beta-1/2)\ge1/4$ follows from
$\beta\ge\sigma>3/4$, while $\delta/2\le1/16$.
The residue term is bounded by $72T^3e^{-T}$ using $X,Y\le T$ and Gamma
exponential decay at $1-\rho$. Thus $\mathcal D\to0$ uniformly over these
zeros. Since $\mathcal D=e^{-1/Y}+\mathcal T$ and $Y\to\infty$,
choose the height threshold so that $e^{-1/Y}\ge3/4$ and
$|\mathcal D|\le1/4$. The reverse triangle inequality gives
$|\mathcal T|\ge1/2$.
\end{proof}
\src{Stronger/InitialDetector.lean; Parameters.lean}{detector_eventually_small; detector_error_scale; detectorTail_decomposition; detectorTail_eventually_large}

\subsection{An exact smooth dyadic partition and finite extraction}
Define
\[
 h(x)=1-S_0(4x-5),\qquad\psi(x)=h(x)-h(2x).
\]
The function $h$ is antitone, equal to one for $x\le5/4$ and zero for
$x\ge3/2$. Thus $0\le\psi\le1$, $\psi$ is smooth, and its support lies
in $[5/8,3/2]\subset(1/2,2)$. For $x>0$,
\[
 \sum_{j\in\Z}\psi(x/2^j)=1.
\]
This is the telescoping sum of $h(x/2^j)-h(x/2^{j-1})$; the two limit
values are one and zero. For a fixed $x$, only finitely many terms are nonzero.
Set
\[
 B_j(\rho)=\sum_{n\ge2}b_X(n)n^{-\rho}e^{-n/Y}\psi(n/2^j).
\]
Absolute summability of the double array follows from nonnegativity of $\psi$
and its partition identity; therefore
$\sum_jB_j(\rho)=\mathcal T(\rho;X,Y)$ and $\sum_j|B_j(\rho)|<\infty$.
Small blocks vanish for $2^j\le X/2$ by M\"obius cancellation and support;
blocks with $j\le0$ also vanish because the $n=1$ term was removed.
\src{Stronger/SmoothPartition.lean; SmoothExpansion.lean; SmoothDetector.lean}{dyadicStep; smoothDyadicWeight; tsum_smoothDyadicWeight; tsum_smoothDyadicBlock; smoothDetectorBlock_zero_of_small}

Let $R=Y(\log T)^2$ and
\[
 J(R)=\{j\in\Z:0\le j\le\lceil\log R/\log2\rceil,\ 2^j\le R\}.
\]
For $T$ sufficiently large, $\#J(R)\le6\log T$.
The discarded block norms, not merely the norm of their sum, satisfy
\[
 \sum_{2^j>R}|B_j(\rho)|\le2T^{-2}.
\]
The source derives this from the tail beginning at
$K=\lfloor R/2\rfloor_+$:
\[
 \sum_{n\ge K}|b_X(n)n^{-\rho}e^{-n/Y}|
 \le\lfloor X\rfloor_+e^{-K/Y}(1-e^{-1/Y})^{-1}.
\]
Here $|b_X(n)|\le\lfloor X\rfloor_+$ is used, $K\ge R/4$ for $R\ge4$,
$e^{-K/Y}\le T^{-4}$ once $\log T\ge16$, and
$(1-e^{-1/Y})^{-1}\le2Y\le2T$. These give $2T^{-2}$.
Since $|\mathcal T|\ge1/2$, at least $1/4$ of the norm sum remains in
$J(R)$; a finite pigeonhole argument gives a $j\in J(R)$ with
\begin{equation}\label{eq:blockdetect}
 X/2<2^j\le Y(\log T)^2,\qquad |B_j(\rho)|\ge1/(24\log T).
\end{equation}
Both the list and its cardinality are fixed before the zero varies.
\src{Stronger/SmoothDetector.lean}{tsum_norm_detector_tail_le; smoothDetectorBlock_scale_tail_le_two_rpow; smoothScaleIndices; exists_smoothDetectorBlock; smoothDetector_eventually_large}

\section{Short blocks: a fixed family, bounded powers and Type I counting}
\label{sec:short}
\subsection{Removing the real-part dependence by an exact Taylor identity}
For a detected scale $2^j$, put
\[
 a_n^{(j)}=b_X(n)e^{-n/Y}\psi(n/2^j),\qquad
 M_{j,e}=2^{\max(j,0)-1+e}\quad(e=0,1),
\]
where the exponent uses natural truncated subtraction as in Lean. A nonzero
block necessarily has $j>0$, in which case $M_{j,0}=2^{j-1}$ and
$M_{j,1}=2^j$.
The smooth block splits exactly into the two natural supports
$(M_{j,0},2M_{j,0}]$ and $(M_{j,1},2M_{j,1}]$.
Both halves retain the \emph{original} smooth scale $2^j$ in $a_n^{(j)}$.
At least one half has norm at least half of $|B_j|$.
\src{Stronger/SmoothShortBlock.lean; ShortFixedFamily.lean}{smoothShortCoefficient; smoothHalfScale; smoothDetectorBlock_eq_halves; exists_smoothShort_half}

For a general positive natural scale $M$ and coefficients $a_n$ with
$|a_n|\le d(n)$ on $(M,2M]$, define
\[
 P_r(t)=\sum_{M<n\le2M}a_n(2n/M-3)^r e^{-\ii t\log(n/M)},\qquad
 q_r(\beta)=\left(\prod_{l=0}^{r-1}\frac{\beta+l}{1+l}\right)(-1/3)^r.
\]
When $0\le\beta\le1$,
$|q_r(\beta)|\le3^{-r}$ and $\sum_r|q_r(\beta)|\le3/2$.
The generalized binomial series gives the exact identity
\begin{equation}\label{eq:tayloridentity}
 \sum_{r\ge0}q_r(\beta)P_r(t)
 = (3/2)^\beta\sum_{M<n\le2M}a_n(n/M)^{-\beta}
                                      e^{-\ii t\log(n/M)}.
\end{equation}
Indeed with $z=2n/M-3\in[-1,1]$, $(2n/(3M))^{-\beta}$ has the displayed
binomial expansion in $z/3$. The finite $n$ sum and absolutely convergent
$r$ sum may be interchanged. Rescaling $n^{-\rho}=M^{-\rho}(n/M)^{-\rho}$
shows that the norm of the right-hand normalized block is
$(3/2)^\beta M^\beta|\sum a_nn^{-\rho}|$.
The coefficients of $P_r$ depend on $a,M,r$, not on $\rho$.
\src{DetectorTaylorFamily.lean; DetectorBinomial.lean; Stronger/GenericTaylorFamily.lean}{detectorTaylorRatio; detectorTaylorWeight; hasSum_detector_binomial; hasSum_genericTaylorPolynomial; genericBlock_eq_normalized}

The triangle bound $|P_r(t)|\le2M^2$ and the geometric weights imply
\[
 \left|\sum_{r\ge J}q_r(\beta)P_r(t)\right|\le3M^2 3^{-J}.
\]
Take $J_T=\lceil10\log T\rceil$. If $M\le T^2$, the tail beyond $J_T$
is at most $3T^{-1}$ because $3^{-J_T}\le T^{-5}$.
For a half-block of norm at least $1/(48\log T)$, this tail is eventually
small enough that a finite Taylor component has norm at least
\begin{equation}\label{eq:taylorheight}
 |P_r(\gamma)|\ge\frac{M^\beta}{192\log T}
          \ge\frac{M^\sigma}{192\log T},\qquad0\le r\le J_T.
\end{equation}
The finite extraction does \emph{not} cost $J_T$: the sum of absolute scalar
weights is at most $3/2$, so some component has at least half the norm of the
finite weighted sum. The scale/half/Taylor index set has cardinality at most
$144(\log T)^2$.
\src{DetectorFamilyScales.lean; Stronger/GenericTaylorFamily.lean; ShortTaylorCover.lean}{detectorTaylorCutoff; norm_fixedGenericTaylor_tail_le; exists_genericTaylor_component; smoothTaylorIndices_card_le_log; smooth_short_Taylor_cover}

\subsection{Bounded power selection and exact convolution supports}
If \eqref{eq:blockdetect} has $2^j>Y^{1/2}$, it is assigned to the long
alternative. Otherwise the selected half-scale satisfies
\[
 T^\delta/4\le M\le T^{c/2}=T^{3a/4}.
\]
For all sufficiently large $T$, there is an integer
$1\le k\le K=\lceil2a/\delta\rceil$ with
\begin{equation}\label{eq:powerlength}
 T^a\le M^k\le T^{3a/2}=T^c.
\end{equation}
To prove this, absorb the factor $1/4$ into $T^{\delta/2}$, set
$x=\log M/\log T$ and choose $k=\lceil a/x\rceil$.
If $x\le a/2$, then $kx<a+x\le3a/2$.
If $a/2<x\le3a/4$, then $k\le2$ and $kx\le3a/2$.
The lower bound $kx\ge a$ and the bound
$k\le\lceil2a/\delta\rceil$ follow from the ceiling inequalities.
Thus the power bound is fixed before every short scale and every zero.
\src{Stronger/PowerScales.lean}{bounded_power_exponent; bounded_power_scale; short_scale_bounded_power; smoothPowerBound}

Restrict the Taylor coefficients to $(M,2M]$ as an arithmetic function
$A_{j,e,r}(n)$, zero elsewhere. Take its genuine Dirichlet convolution power
$A_{j,e,r}^{*k}$, not a pointwise power of coefficients.
The product Dirichlet series is the $k$th power of the finite block series,
with support in $(M^k,(2M)^k]$. The source proves this finite convolution
identity and its support before using it. Divide that support into the
$k$ dyadic intervals $(2^lM^k,2^{l+1}M^k]$, $0\le l<k$.
At least one piece has norm at least $|P_r(\gamma)|^k/k$.
Use the coefficientwise complex conjugate of the convolution power to change
the negative phase convention to the positive detecting polynomial while
preserving its norm at the \emph{same} ordinate $\gamma$.
Define
\[
 L_{j,e,r,k,l}=2^lM_{j,e}^k,\qquad
 a_{j,e,r,k,l}(n)=T^{-\eta}\overline{A_{j,e,r}^{*k}(n)}.
\]
The source's bounded-convolution divisor estimate makes these coefficients
uniformly at most one whenever their entire support is below $T^2$, for fixed
$K$ and arbitrary positive $\eta$. Equations \eqref{eq:powerlength} and
$c<1$ ensure that support condition eventually, since
$(2M)^k\le2^KT^c\le T$.
These estimates preserve a bound $N^\eta$ because, for an arithmetic
function with $|A(n)|\le d(n)$, the source first proves
$|A^{*k}(n)|\le d(n)^{2k}$. Use the divisor subpower estimate with loss
$\eta/(2(K+1))$ and $k\le K$ to obtain
$|A^{*k}(n)|\le C^{2(K+1)}n^\eta$. To normalize on $n\le T^2$, use that
estimate with loss $\eta/4$ and then absorb its fixed constant into
$T^{\eta/2}$; the resulting bound is $T^\eta$ before any family member
is selected.
The height loss is at most
$C_K(\log T)^K$, with
$C_K=2^K(K+1)192^K$; it is eventually at most $T^\eta$.
After normalization the selected piece therefore satisfies
\begin{equation}\label{eq:poweredheight}
 |D_{L_q,a_q}(\gamma)|\ge L_q^\sigma T^{-2\eta},\qquad
 T^a\le L_q\le2^KT^c.
\end{equation}
\src{Stronger/FiniteConvolution.lean; ShortFixedFamily.lean; ShortPoweredHeight.lean}{exists_finiteBlock_powered_piece; smoothPoweredCoefficient; normalizedSmoothPoweredCoefficient_eventually; exists_normalizedSmoothPowered_height}

The \emph{fixed family} consists of all tuples
$q=(j,e,r,k,l)$ with $j\in J(Y\log^2T)$, $e=0,1$, $0\le r\le J_T$,
$1\le k\le K$, $0\le l<k$, and the two conditions in
\eqref{eq:powerlength} for that half-scale. These restrictions depend on
$\sigma,\delta,T$, not on any zero. Its cardinality is at most
$144K^2(\log T)^2$.
Only the chosen index depends on the zero; the family coefficients never do.
Every slab zero has either the long block or a member of this family satisfying
\eqref{eq:poweredheight} at its original ordinate.
\src{Stronger/ShortPoweredFamily.lean}{smoothPoweredIndices; smoothPoweredFamilyCoefficient; smoothPoweredIndices_card_eventually; smooth_long_or_short_powered_cover}

\subsection{Jensen multiplicity restoration and the short count}
\begin{lemma}[Uniform local multiplicity]\label{lem:localzeros}
There is $C_b>0$ such that for every $\sigma\ge3/4$, $T\ge2$ and integer $k$,
\[
 \sum_{\rho\in\mathcal Z(\sigma;T,2T)\atop k\le\Ima\rho<k+1}m(\rho)
                  \le C_b\log(2T).
\]
\end{lemma}
\begin{proof}
For $T\ge8$ and a nonempty bin, its center is
$z_k=2+\ii(k+1/2)$; the bin location implies $T-1\le k\le2T$.
Each relevant zero has $\Rea\rho\ge7/10$ and lies inside radius $3/2$ about
$z_k$, because its horizontal difference is at most $13/10$ and its vertical
difference at most $1/2$.
The disk of radius $7/4$ avoids $1$. The Abel zeta growth estimate on its
sphere gives $|\zeta(s)|\le100T^3$.
On the center line $\Rea s=2$, the absolutely convergent M\"obius inverse
series has norm less than $5/3$, since $\zeta(2)=\pi^2/6<5/3$.
Its product with zeta is one, so $|\zeta(z_k)|\ge3/5=0.6$.
Jensen's formula for radii $3/2,7/4$, and the source's finite-sum analytic-order
to-divisor bridge, give
\[
 \sum_{\text{bin}}m(\rho)
 \le\frac{\log(100T^3/0.6)}{\log((7/4)/(3/2))}\ll\log(2T).
\]
The bridge is necessary: it identifies each finite natural vanishing order
with the nonnegative analytic divisor, rather than counting distinct roots.
For $2\le T\le T_0$, all bin sums are bounded by the finite count
$N(3/4,2T_0)$. Since $\log(2T)\ge\log4>0$, this bounded contribution is
absorbed by increasing the same $C_b$.
\end{proof}
\src{ZetaJensenInputs.lean; LocalZeroCount.lean}{euler_product_lower_bound_2; finset_analyticOrderNatAt_le_finsum_divisor; zeta_unit_bin_sum_le_jensen; zeta_local_multiplicity_bound}

For any finite subset $S$ of these zeros, select the parity class with at
least half the occupied unit bins, then one zero in each selected bin.
Distinct selected bins of the same parity differ by at least two; their
ordinates are one-separated. The ordinate map on the selected representatives
is injective. Lemma~\ref{lem:localzeros} implies
\begin{equation}\label{eq:weightedextract}
 \sum_{\rho\in S}m(\rho)\le2C_b\log(2T)\,|U|,
\end{equation}
where $U$ is their ordinate set. This is a weighted finite extraction, retaining
all the lost multiplicities in the factor $C_b\log(2T)$.
\src{ZeroSeparation.lean}{ordinateBin; occupied_bins_parity; weighted_separated_extraction}

Let $S_{\rm short}$ be exactly those zeros in $\mathcal Z(\sigma;T,2T)$
detected by \emph{some} fixed-family member at level
$V_q=L_q^\sigma T^{-2\eta}$; set
$N_{\rm short}=\sum_{\rho\in S_{\rm short}}m(\rho)$.
Choose $\eta>0$ with $8\eta<\sigma-3/4$ and
$\kappa=(\sigma-3/4)/2>0$.
Use the comparison ambient height
\[
 H_T=32T(\log T)^2,
\]
which eventually satisfies $T\le H_T\le T^2$ and $H_T^{1/3}\le\sqrt T$.
The fixed scales obey $L_q\ge T^a\ge\sqrt T$, $L_q\le T\le H_T$ and
\[
 V_q=L_q^\sigma T^{-2\eta}\ge L_q^{\sigma-4\eta}
                                     \ge L_q^{3/4+\kappa}.
\]
Thus each family member satisfies \emph{all} hypotheses of Theorem~\ref{thm:LV}
with height $H_T$; the slab is contained in an interval of that length.
Partition the separated ordinate representatives into the superlevel subsets
for each family member and use the finite union-cardinality bound.
Combining with \eqref{eq:weightedextract} gives
\begin{equation}\label{eq:shortcover}
 N_{\rm short}\le C\log(2T)H_T^\eta
       \sum_q\left(\frac{L_q^2}{V_q^2}+\frac{L_q^3\sqrt{H_T}}{V_q^4}\right).
\end{equation}
Overlapping detection sets only enlarge this upper bound.
\src{Stronger/ShortZeroCount.lean; ShortExponentAccounting.lean}{smoothShortZeroFinset; zetaSubset_large_values_cover_indexed; smoothPoweredFamily_largeValues_admissible}

The two summands are, exactly,
\[
 L_q^{2-2\sigma}T^{4\eta},\qquad
 L_q^{3-4\sigma}\sqrt{H_T}\,T^{8\eta}.
\]
The first length exponent is nonnegative; use the upper length
$2^KT^c$. The second is negative; use the lower length $T^a$.
Equations~\eqref{eq:exponentidentities}, and
$\sqrt{H_T}\le32T^{1/2}\log T$, give a bound by
$(2^K+32)\log T\,T^{p+8\eta}$ per member.
The family contributes at most $144K^2\log^2T$; the bin factor at most
$2\log T$; $H_T^\eta\le T^{2\eta}$; and the remaining four logarithms
are eventually at most $T^{2\eta}$. Therefore
\begin{equation}\label{eq:shortcount}
 N_{\rm short}\le C_{\sigma,\delta,\eta}T^{p+12\eta}.
\end{equation}
This accounts for $8\eta$ in the quotient terms, $2\eta$ in the ambient
height and $2\eta$ in the logarithms. The fixed $2^K$ enlargement is not
omitted, and no long-zero estimate was used.
\src{Stronger/ShortExponentAccounting.lean; ShortZeroCount.lean}{smooth_largeValue_term_le; smooth_family_sum_eventually; smoothShortZeroCount_bound}

\section{Long blocks, weighted sampling and the full density endpoint}
\label{sec:long}
Define the complementary zero set
$S_{\rm long}=\mathcal Z(\sigma;T,2T)\setminus S_{\rm short}$ and its
multiplicity count $N_{\rm long}$. This is an actual complement, not an
independently selected family of difficult zeros. Hence
\begin{equation}\label{eq:partition}
 N_{\rm slab}(\sigma;T,2T)=N_{\rm short}+N_{\rm long}.
\end{equation}
The fixed-family cover implies that every member of $S_{\rm long}$ has a
large smooth block with $2^j>Y^{1/2}$.

\subsection{The uniform smooth Mellin identity and its residue}
For $q\ge0$ put $\psi_q(x)=\psi(x)e^{-qx}$ and
$\widehat\psi_q(s)=\int_0^\infty\psi_q(x)x^{s-1}\dd x$.
The compact support lies in $[5/8,3/2]$, so this Mellin transform is entire.
Its rapid vertical decay is uniform in \emph{all} $q\ge0$ on any fixed
real strip. For each derivative order, the inequality
$q^ne^{-qx}\le n!/(5/8)^n$ on the support gives uniform derivative bounds;
Leibniz differentiation and Fourier decay of the logarithmic lift give
$|\widehat\psi_q(u+\ii v)|\le C_A(1+|v|)^{-A}$ and uniformly weighted
$L^1$ bounds. Thus the damping parameter $q=2^j/Y$ need not be bounded.
\src{Stronger/DampedCutoff.lean; UniformMellin.lean; MellinEntire.lean}{dampedCutoff_uniform_derivatives; dampedCutoff_uniform_mellin_decay; dampedCutoff_uniform_mellin_weighted_L1; dampedCutoff_mellin_entire}

For $N>0$ the smooth mollifier block is
\[
 B(\rho;X,N,q)=\sum_{n\ge1}b_X(n)n^{-\rho}\psi(n/N)e^{-qn/N}.
\]
It equals the normalized integral of
$N^s\widehat\psi_q(s)\zeta(\rho+s)M_X(\rho+s)$ on $\Rea s=1/2$.
Shift to $\Rea s=1/2-\beta$. The only pole is $s=1-\rho$ and its residue is
\[
 R_B=N^{1-\rho}\widehat\psi_q(1-\rho)M_X(1).
\]
There is no Gamma pole at zero in this smooth Mellin kernel; in particular the
block-contour identity does not require the zero hypothesis. The smooth
Mellin bounds and Abel zeta growth prove both-line integrability and vanishing
horizontal edges. For $T\le|\gamma|$, $0\le X\le T$,
$1\le N\le T^2$, the residue is at most $T^{-100}$ eventually,
uniformly in $\beta\in[3/4,1)$ and $q\ge0$; a sufficiently high fixed
Mellin-decay power absorbs all powers of $X,N,T$.
For $N=2^j\ge2$ the removed $n=1$ term has zero smooth weight, so this is
exactly $B_j(\rho)$ with $q=N/Y$.
\src{Stronger/SmoothMellinIdentity.lean; SmoothContourIdentity.lean; SmallResidue.lean; SmoothDetector.lean}{smoothMollifierBlock_eq_rightContour; smoothMollifierBlock_contour_identity; smoothDetectorResidue_eventually_small; smoothDetectorBlock_eq_smoothMollifierBlock}

Define, for integer $A\ge3$,
\[
 w_A(u)=(1+|u|)^{-A},\qquad
 \mathcal M_A(\gamma)=\int_{\R}w_A(u)\zabs(\gamma+u)\dd u.
\]
The uniform left-line Mellin bound and $|M_X(s(\gamma+u))|\le2\sqrt X$
give
\[
 |B(\rho;X,N,q)|\le|R_B|+C\sqrt XN^{1/2-\beta}\mathcal M_A(\gamma).
\]
The residue is at most $1/(48\log T)$ eventually. Hence a block of norm at
least $1/(24\log T)$ satisfies
$1/(48\log T)\le C\sqrt XN^{1/2-\beta}\mathcal M_A(\gamma)$.
For $N>Y^{1/2}$ and $\beta\ge\sigma$, the negative exponent gives
$N^{1/2-\beta}\le Y^{-(\sigma-1/2)/2}$, so every long zero satisfies
\begin{equation}\label{eq:longmean}
 Y^{(\sigma-1/2)/2}\le48C\sqrt X\log T\,\mathcal M_{128}(\gamma).
\end{equation}
We have selected the source's fixed decay order $A=128$ at this point.
\src{Stronger/WeightedCriticalMean.lean; LongBlockDetection.lean; LongMeanCounting.lean}{weightedCriticalMean; uniform_smoothMollifierBlock_mean_bound; smoothMollifierBlock_eventually_detects_mean; smooth_long_mean_or_short_powered_cover}

\subsection{Twelfth-power sampling, including the full-line tail}
For $A\ge2$, $0<\int w_A\le\pi$ and a nonnegative integrable weighted
function $f$ satisfies
\[
 \left(\int w_A f\right)^{12}\le\left(\int w_A\right)^{11}
                                      \int w_A f^{12}.
\]
The source proves this by the tangent inequality
$12a^{11}x\le x^{12}+11a^{12}$ for $a,x\ge0$: take
$a=(\int w_Af)/(\int w_A)$, multiply by $w_A$, integrate and rearrange.
All three required weighted integrals are proved integrable first. For
$A=128$, coarse growth supplies the twelfth weighted integral.
For a one-separated finite $U$ the kernel row sum
$\sum_{u\in U}w_{128}(t-u)\le S$ is uniform, by comparison with
$(1+|t-u|)^{-2}$ and the earlier separation bound.

On the complement of $[0,3T]$, for $u\in[T,2T]$,
$|t-u|\ge T$ and $|t|\le3|t-u|$. Thus
\begin{equation}\label{eq:physicaltail}
 w_{128}(t-u)\zabs(t)^{12}
       \le12^{12}T^{-114}w_2(t-u).
\end{equation}
To verify it, coarse growth gives
$\zabs(t)\le12(1+|t-u|)$; cancel twelve of the 128 denominator powers,
reserve two powers for integrability, and use the remaining 114 powers to
produce $T^{-114}$. Integrating gives at most
$12^{12}\pi T^{-114}$ per sample. One-separated $U\subset[T,2T]$ has
$|U|\le T+1\le2T$, so the total sample tail costs at most
$2\cdot12^{12}\pi T^{-113}$.
Consequently
\begin{equation}\label{eq:sampling}
 \sum_{u\in U}\mathcal M_{128}(u)^{12}
 \le\pi^{11}S\int_{[0,3T]}\zabs(t)^{12}\dd t
            +2\cdot12^{12}\pi^{12}T^{-113}.
\end{equation}
The whole-line means have not been replaced by a compact integral without
accounting for the tails.
\src{Stronger/MomentPower.lean; SeparatedMoment.lean; PhysicalTwelfthTail.lean; SeparatedCriticalSampling.lean}{integral_weighted_twelfth; momentDecay_mass_bounds; weightedCriticalTwelfth_tail_pointwise; separated_weightedCriticalMean_twelfth_le}

Select weighted separated representatives of $S_{\rm long}$ using
\eqref{eq:weightedextract}, preserving the property \eqref{eq:longmean}.
Raise \eqref{eq:longmean} to the twelfth power and sum over their ordinate set:
\[
 |U|Y^{6\sigma-3}\le(48C)^{12}X^6(\log T)^{12}
                              \sum_{u\in U}\mathcal M_{128}(u)^{12}.
\]
Here $(\sqrt X)^{12}=X^6$ and
$(Y^{(\sigma-1/2)/2})^{12}=Y^{6\sigma-3}$ exactly.
Use \eqref{eq:sampling} and the local multiplicity factor
$\log(2T)\le2\log T$. The result is the unconditional reduction
\begin{equation}\label{eq:longraw}
 N_{\rm long}Y^{6\sigma-3}
 \le C X^6(\log T)^{13}
          \left(\int_{[0,3T]}\zabs(t)^{12}\dd t+T^{-113}\right).
\end{equation}
Its right side contains the actual moment, not an assumed upper bound on it.
\src{Stronger/LongMeanCounting.lean; LongMultiplicityReduction.lean}{long_mean_twelfth_count; long_mean_count_le_physical_moment; smoothLongZeroCount_physical_moment_reduction}

\subsection{Exact long saving and final parameter choice}
Use Theorem~\ref{thm:moment} with a positive loss
$\nu\le(1-c)/8$. If its constant is $D>0$, then for $T\ge1$ eventually,
\[
 \int_{[0,3T]}\zabs^{12}+T^{-113}\le(D+1)T^{2+\nu}.
\]
Dividing \eqref{eq:longraw} by the positive $Y^{6\sigma-3}$ gives the scalar
factor
\begin{equation}\label{eq:longscalar}
 \frac{X^6(\log T)^{13}T^{2+\nu}}{Y^{6\sigma-3}}
   =(\log T)^{13}T^{p-(1-c)+6\delta+\nu}.
\end{equation}
Since $6\delta<(1-c)/2$, the pre-logarithmic exponent is less than
$p-3(1-c)/8$. Absorb $(\log T)^{13}$ using the further allowance
$(1-c)/8$ to obtain
\begin{equation}\label{eq:longcount}
 N_{\rm long}\le C' T^{p-(1-c)/4}
\end{equation}
eventually. This fixed saving is possible precisely because $c<1$.
\src{Stronger/LongScalarBudget.lean; ConditionalDensity.lean}{long_count_scale_identity; long_count_scalar_log_eventually; smoothLongZeroCount_bound_of_twelfth_moment}

\begin{proof}[Proof of Theorem~\ref{thm:main}]
First suppose $3/4<\sigma<1$. Given $\varepsilon>0$, choose, in the same
order as the source,
\[
 \delta=\min(1/16,(1-c)/24),\qquad
 \eta=\min(\varepsilon/24,(\sigma-3/4)/16),\qquad
 \nu=(1-c)/16.
\]
They satisfy $\delta>0$, $\delta\le1/8$, $\delta<(1-c)/12$,
$\eta>0$, $8\eta<\sigma-3/4$, $12\eta\le\varepsilon$,
and $0<\nu\le(1-c)/8$. Choose the actual moment constant and threshold
with this $\nu$. Thus every parameter, power bound and coefficient family
is fixed before heights and zeros vary.

Combine the exact partition \eqref{eq:partition} with
\eqref{eq:shortcount} and~\eqref{eq:longcount}. Since $T\ge1$ and $1-c>0$,
the long power is at most $T^{p+12\eta}$. Therefore
\[
 N_{\rm slab}(\sigma;T,2T)\le C_s T^{r},\qquad r=p+12\eta>0,
\]
eventually in $T$.

We detail the source's conversion from eventual slabs to a bound at every
height, because dropping the compact part would weaken the stated endpoint.
Let $R\ge1$ exceed the slab threshold and
$Q=N_{\rm slab}(\sigma;0,2R)$. If $1\le T<R$, monotonicity gives
$N_{\rm slab}(\sigma;T,2T)\le Q\le QT^r$; if $T\ge R$, use the eventual
bound. Hence the same estimate, with constant $C_s+Q$, holds for every
$T\ge1$.
Put $b=2^r>1$, $Q_0=N_{\rm slab}(\sigma;0,1)$ and
$A=Q_0+(C_s+Q)/(b-1)>0$. Then $A+(C_s+Q)\le Ab$.
Inductively split the inclusive interval $[0,2^{k+1}]$ at $2^k$ to obtain
\[
 N_{\rm slab}(\sigma;0,2^k)\le A(2^k)^r.
\]
The shared boundary may be counted twice by this inequality, which is harmless
for an upper bound and preserves every zero. For arbitrary $T\ge1$, choose
$2^k\le T<2^{k+1}\le2T$. Monotonicity yields
$N_{\rm slab}(\sigma;0,T)\le AbT^r$.
Conjugation and its multiplicity identity give
$N(\sigma,T)\le2AbT^r$.
Since $r=p+12\eta\le p+\varepsilon$ and $T\ge2$, enlarge the power to
$T^{p+\varepsilon}$. The resulting constant is positive and chosen before
all $T$.

For $\sigma\ge1$ the zero set is empty by Section~\ref{sec:zeros}.
Take $C=1$; the right side is positive for $T\ge2$, even if its exponent
is negative. Finally use the exact membership lemma to identify
$\mathcal Z(\sigma;-T,T)$ with the finite set of \emph{all} zeros in the
registered statement. Lemma~\ref{lem:zeroorder} supplies positive finite
multiplicity and the equality of natural and extended order at every member.
The source's natural-to-real sum casts merely change the numerical encoding
of the same sum. This proves the complete endpoint, including its witnesses
and all boundary conventions.
\end{proof}
\src{Stronger/FinalParameters.lean; ConditionalDensity.lean; NativeDensity.lean; DensityDyadicSummation.lean; Solution.lean}{exists_stronger_loss_parameters; stronger_slab_bound_of_twelfth_moment; zetaDensity_bound_of_eventual_slabs; stronger_density_bound_native; DensityStronger.density_bound}

\section{Reading guide, interfaces and fidelity checks}
\label{sec:fidelity}
A natural learning order is: the actual-zero definitions in
Section~\ref{sec:zeros}; the Type I theorem in Section~\ref{sec:typeI}; the
fourth-moment source and point-mean inequality in
Sections~\ref{sec:fourth}--\ref{sec:pointmean}; the signed stationary source
and its Gram packet in Sections~\ref{sec:atkinson}--\ref{sec:packets}; the
peak occupancy and high/low moment calculus in Section~\ref{sec:peaks};
then the smooth detector, fixed short family and long sampling in
Sections~\ref{sec:detector}--\ref{sec:long}.
The moment and Type I chains can be studied independently until the density
consumer is assembled.

\paragraph{Conditional interfaces are not assumptions of the conclusion.}
The names \lean{PointMeanInput}, \lean{LocalMeanPacketInput},
\lean{LocalMeanStationaryInput}, \lean{LocalMeanSourceToPrefix} and the
literal-twelfth-moment premise describe reusable \emph{consumer} propositions.
Their mathematical statements are given above. The actual point-mean producer
fills the first; the actual local stationary producer and signed-prefix adapter
fill the next three; the native high/low moment fills the last. In the final
\lean{NativeDensity} composition no one of these remains a hypothesis.
Some older source comments call these interfaces ``unproved'' within their own
conditional module; the native composition, not that isolated comment, is the
relevant dependency chain. This blueprint does not add these interfaces to the
assumptions of the registered endpoints.

\paragraph{Endpoint-sensitive features retained.}
\begin{enumerate}
\item The proof counts analytic orders, not just distinct ordinates; finite
unit-bin extraction is followed by logarithmic multiplicity restoration.
\item M\"obius cancellation removes the $n=1$ detector term explicitly.
The smooth partition and tail estimates are for the resulting $n\ge2$ series.
\item Smooth dyadic support is $[5/8,3/2]$; half-block splitting keeps the
original $2^j$ scale in both coefficient functions.
\item The Taylor family coefficients are independent of the zero's real part.
They are selected from a fixed finite family; its chosen index, not its
coefficients, may depend on the zero.
\item Convolution is genuine Dirichlet convolution; the power $k$ is bounded
before $T$, and the enlarged support keeps the factor $2^K$.
\item Type I comparison is used with $32T\log^2T$, not with an unchecked
scale. Both signs of the length exponents are used in the proper direction.
\item The Atkinson integer-lattice phase $e^{-2\pi\ii x}$ is retained in the
continuous integral. Both Bessel branches are evaluated or bounded, with their
whole series and normalizations tracked.
\item The negative stationary saddle has its own residual weight and the
$-\ii$ Fresnel factor. It is not silently merged with the positive saddle.
\item Maximal prefixes are height-dependent, and pairwise Gram entries use
the minimum of the two prefixes. The far-gap triangle branch is retained.
\item The smooth long-contour residue and the full-line weighted-twelfth tail
are bounded explicitly. The tail exponents $-114$ and $-113$ enter before
the final physical-moment bound.
\item The short zero set and its complement partition the original slab.
The final dyadic summation includes the compact initial count and closed
height boundaries, then uses conjugation.
\item No claim of $\varepsilon=0$, $\sigma=3/4$, effective constants or
boundary-uniformity is inserted.
\end{enumerate}

\paragraph{Source attribution and verification boundary.}
The repository credits Mark Lewko for the development and selected analytic
proofs adapted from Scott McColm, Conor Grogan and other credited sources.
The local source retains its AGPL-3.0-only and third-party notices; those notices
are preserved with the archived snapshot. This mathematical reconstruction is
not a new Lean certificate or a fresh independent kernel replay. The recorded
registry verification applies to its pinned source and toolchain. The source
references here support an audit of the reconstruction against that source;
LaTeX compilation verifies typesetting, not mathematical truth. The auxiliary
\lean{Challenge.lean} contains specification holes and is not imported by the
production solution; no such hole is used as an informal proof step.

\begin{thebibliography}{9}
\raggedright
\bibitem{registry}
Palomar Registry, \emph{A zero-density estimate with exponent
$3(1-\sigma)/(2\sigma)$ for the Riemann zeta function},
PALOMAR-2026-10-09-000005, version 1.
\url{https://data.palomar-registry.org/entries/PALOMAR-2026-10-09-000005-v1.json}.
\bibitem{source}
Mark Lewko and credited contributors, \emph{Registered Lean development},
revision \texttt{804bdb55b1fc40e6a1a6878fd96078644a70e288}.
\url{https://github.com/mrklwk/zeta-density-three-over-two-sigma-lean/tree/804bdb55b1fc40e6a1a6878fd96078644a70e288}.
\bibitem{mathlib}
The mathlib community, \emph{mathlib4}, source revision
\texttt{065356127b1dc0016f66b7283ce0ce2c4055aa55}.
\url{https://github.com/leanprover-community/mathlib4/tree/065356127b1dc0016f66b7283ce0ce2c4055aa55}.
\bibitem{hb}
D.~R. Heath-Brown, 1978 twelfth-moment source cited by the registered
formalisation, especially Lemma 3 and equation (44).
\url{https://wiki.math.ntnu.no/_media/ma3001/2025h/analyticnumbertheory/heathbrowntwelfthmoment.pdf}.
This reference is attribution and context; the proofs reconstructed here follow
the Lean sources rather than replacing them with the paper's argument.
\end{thebibliography}
\appendix
\section{Analytic adapters behind the native producers}
\label{app:adapters}
This appendix expands selected analytic adapters whose exact construction is
important for reading the main proof. They are locally proved results, not
extra assumptions. All branches and finite-limit conventions here agree with
the production imports.

\subsection{The sharp approximate functional equation used for Weyl growth}
For $s\in\C$ and a real cutoff $x$, let
$S(s,x)=\sum_{1\le n\le\lfloor x\rfloor_+}n^{-s}$ and define the literal
remainder
\[
 R_{\rm AFE}(s;x,y)=\zeta(s)-S(s,x)-\chi(s)S(1-s,y),
 \quad \chi(s)=2^s\pi^{s-1}\Gamma(1-s)\sin(\pi s/2).
\]
The unit term is included in both sums. Defining this remainder alone supplies
no estimate on it; the native estimate is the following.

For $t\ge1$, $s=1/2+\ii t$, assume
\begin{gather*}
 x,y\ge1,\quad y\le x,\quad 2\pi xy=t,\quad x=k+1/2\ (k\in\Z),\\
 \lfloor y\rfloor_++1/8\le y\le\lfloor y\rfloor_++1/2,
 \qquad e^{-\pi t}(t^2+1)\le1.
\end{gather*}
Then $|R_{\rm AFE}(s;x,y)|\le C_{\rm AFE}$, where the source gives the
explicit positive constant
\begin{align*}
 C_P&=\frac{\log2+1}{2\pi}
 +\frac{1+2\pi}{4\pi^2}(\pi/2+\log2)
 +\frac{2(1+2\pi)}{4\pi^3}\,9
 +\frac{1+2\pi}{4\pi^3}\,17,\\
 C_{\rm AFE}&=C_P+3+24/\pi+(1-e^{-\pi})^{-1}.
\end{align*}
Here is the source's decomposition and error proof. Define the actual finite
oscillatory integral
\[
 J(s;a,b;m)=\int_a^b u^{-s}e^{2\pi\ii mu}\dd u
\]
and $J_\infty(s;a;m)$ as the limit as $b\to\infty$ of this integral for
$0<\Rea s<1$, $m>0$. It is not an undamped whole-half-line Lebesgue
integral. Let $M=\lfloor y\rfloor_+$ and
\begin{gather*}
 B=\frac{x^{1-s}}{s-1},\qquad F=\sum_{m=1}^M J_\infty(s;x;m),\qquad
 F_0=\sum_{m=1}^M J_\infty(s;0;m),\\
 L=\sum_{m=1}^M J(s;0,x;m),\qquad Q=\chi(s)S(1-s,y).
\end{gather*}
The Abel sawtooth identity for actual zeta, followed by its finite Poisson
expansion and the proved limit of its positive and negative frequency tails,
gives
\[
 |\zeta(s)-S(s,x)-B-F|\le C_P.
\]
For clarity, the four contributions to $C_P$ are the endpoint first-order
coefficient, the paired harmonic endpoint coefficient, and the two higher
coefficient tails. On the half-integral $x$, alternating harmonic tails have
their explicit bounds. The two distance-to-integer gaps in $y$ give
\[
 \text{minus-square tail}+\text{plus-square tail}\le9/y,
 \qquad
 \text{minus-cube tail}+\text{plus-cube tail}\le17/y.
\]
The differentiated endpoint amplitudes satisfy
$H/y\le1+2\pi$, $H_1/y\le2(1+2\pi)$ and
$H(y/x)/y\le1+2\pi$, because $y\le x$ and $x\ge1$.
Substitution into the four terms yields exactly $C_P$.

The elementary pole term obeys $|B|\le1$.
For the lower integrals, integration by parts in their logarithmic phase and
retention of the half-integer endpoint produce the arithmetic gap sum
\[
 \sum_{m=1}^M\frac{m^2}{y^2(y-m)}\le8(1+\log y).
\]
The gap $y-M\ge1/8$ controls the closest term; for the other terms the
harmonic sum comparison gives the logarithm. Since
$x^{-1/2}(1+\log y)\le3$ for $1\le y\le x$, and the two boundary factors
are each at most one, this proves $|L|\le2+24/\pi$.

To evaluate $F_0$, introduce positive damping $e^{-\epsilon u}$.
The convergent Gamma--Laplace integral is
\[
 \int_0^\infty u^{-s}e^{2\pi\ii mu-\epsilon u}\dd u
   =(\epsilon-2\pi\ii m)^{s-1}\Gamma(1-s),\qquad\epsilon>0.
\]
On any finite interval, dominated convergence removes the damping. On an upper
interval starting at $R>t/(\pi m)$, the phase derivative has magnitude at
least $\pi m$, and the nonnegative amplitude $u^{-\Rea s}e^{-\epsilon u}$
is decreasing. The proved primitive bound and integration by parts give the
uniform tail bound $2R^{-\Rea s}/(\pi m)$. Thus the finite-interval Abel
limit and this tail bound identify
\[
 J_\infty(s;0;m)=(-2\pi\ii m)^{s-1}\Gamma(1-s).
\]
The negative-imaginary-axis principal branch is retained. Its exact factorization
is
\[
 \Gamma(1-s)(2\pi/\ii)^{s-1}
       =\chi(s)(1+e_\Gamma(s)),\qquad
 e_\Gamma(s)=\frac{e^{\pi\ii s}}{1-e^{\pi\ii s}}.
\]
The denominator has norm at least $1-e^{-\pi|t|}$; hence for $t\ge1$,
$|e_\Gamma(s)|\le e^{-\pi t}/(1-e^{-\pi})$.
Also $|S(1-s,y)|\le y^{1/2}\log y+1\le t^2+1$.
Since $|\chi(s)|=1$, this gives $|F_0-Q|\le(1-e^{-\pi})^{-1}$.
The exact relation $F=F_0-L$ and
\[
 R_{\rm AFE}=(\zeta-S-B-F)+B+(F-Q)
\]
prove the displayed constant bound. This makes all three analytic error sources
explicit, rather than postulating a sharp functional equation.

Finally the source constructs the cutoffs at every $t\ge200\pi$.
Let $r=\sqrt{t/(2\pi)}\ge10$, set
\[
 M=\lfloor r/2\rfloor_+,\quad b=M+1/2,\quad u=r^2/b,\quad
 A=\lfloor u+1/2\rfloor_+,\quad x=A+1/2,\quad y=r^2/x.
\]
The floor inequalities give $u<x\le u+1$, and consequently
\[
 r\le x\le(12/5)r,\quad(2/5)r\le y\le r,\quad
 M+1/8\le y\le M+1/2,\quad\lfloor y\rfloor_+=M.
\]
The first scale range and $2\pi>6$ give $A^2\le t$; the second gives
$M^2\le t$. The exponential envelope is eventually at most one.
These are precisely the short-polynomial cutoffs used in
Proposition~\ref{prop:weyl}.
\src{WeylPort/AFE/Objects.lean; WeylPort/AFE/DampedWeightedKernel.lean; WeylPort/AFE/ChiGammaFactor.lean; WeylPort/AFE/UniformError.lean; WeylPort/CutoffSelection.lean}{afeRemainder; weightedIntegralTail_zero_eq_gamma; gamma_factor_eq_chi_mul_one_add; norm_critical_afe_uniform; exists_afe_cutoffs_all_heights}

\subsection{The literal Neumann kernel: contour, branch and error}
The exact ray proof of \eqref{eq:Ytwo} is short enough to give explicitly.
Define, with principal half-powers,
\[
 K_x(z)=e^{\ii xz}(1-z)^{-1/2}(1+z)^{-1/2},\qquad
 V_x(b)=\int_0^\infty K_x(b+\ii t)\dd t.
\]
For $0\le b<1$ a rectangle in the upper half-plane, its upper edge bounded by
$C/H$ times $e^{-xH}$, yields the equality of the two vertical integrals and
the finite real edge. Under $z=\sin u$, the positive cosines cancel the two
principal half-powers, giving
\[
 \int_0^{\pi/2}e^{\ii x\cos u}\dd u=\ii V_x(0)-\ii V_x(1).
\]
The source takes $b=\sin(\pi/2-H^{-1})\to1$ and proves convergence of the
vertical integral at $b=1$, where the local endpoint majorant is $t^{-1/2}$.
The $b=0$ integral is the real tail
$\int_0^\infty e^{-xt}/\sqrt{1+t^2}\dd t$.
Taking imaginary parts and comparing with the Schl\"afli definition gives
$Y_0(x)=-(2/\pi)\Rea V_x(1)$.

On the unit ray,
\[
 V_x(1)=e^{\ii x}\frac{1+\ii}{2}
       \int_0^\infty t^{-1/2}e^{-xt}(1+\ii t/2)^{-1/2}\dd t.
\]
This factorization uses multiplication of a positive real base with a
nonzero complex base, which does not cross the principal logarithm's cut.
The constant is exactly $(-\ii)^{-1/2}2^{-1/2}=(1+\ii)/2$.
Put $w=(1+\ii u)^{1/2}$. Then
$\Rea w\ge1$, $|w|\ge1$, $|w+1|\ge2$, and
$w-1=\ii u/(w+1)$. The algebraic identity
\[
 2w\bigl(w^{-1}-1+\ii u/2\bigr)=(w-1)^2(w+2)
\]
therefore implies the global bound
\[
 |(1+\ii u)^{-1/2}-1+\ii u/2|\le3u^2/8.
\]
At $u=t/2$, multiply by $t^{-1/2}e^{-xt}$ and integrate.
The three exact moments are
\[
 \int_0^\infty t^{k-1/2}e^{-xt}\dd t
    =\Gamma(k+1/2)x^{-k-1/2},\quad k=0,1,2.
\]
Their values are $\sqrt\pi x^{-1/2}$,
$(\sqrt\pi/2)x^{-3/2}$ and $(3\sqrt\pi/4)x^{-5/2}$.
Thus the Laplace integral differs from
$\sqrt\pi x^{-1/2}-\ii\sqrt\pi x^{-3/2}/8$ by at most
$(9/128)\sqrt\pi x^{-5/2}$.
Multiplication by $e^{\ii x}(1+\ii)/2$, taking its real part, and multiplying
by $-2/\pi$ give \eqref{eq:Ytwo} and its stated error. All these integrals
are absolutely convergent for $x>0$; no asymptotic series or unproved Bessel
property enters the computation.
\src{Stronger/Atkinson/NeumannContourShift.lean; NeumannSchlafliEntry.lean; NeumannLaplaceRepresentation.lean; NeumannRayFactorization.lean; NeumannLaplaceAmplitude.lean; NeumannLaplaceRemainder.lean}{neumann_schlafli_complex_identity; dfiBesselY0_eq_neumannLaplaceIntegral; norm_neumannLaplaceAmplitude_sub_linear_le; norm_neumannLaplaceIntegral_sub_approximation_le}

\subsection{Finite Fresnel windows and the symmetric stationary cancellation}
For $c>0$, let
$F_c(H)=\int_{-H}^H e^{-2\pi\ii cz^2}\dd z$.
For positive damping, the whole-line Gaussian value is
$G_\epsilon(c)=(\pi/(\epsilon+2\pi c\ii))^{1/2}$.
The phase derivative outside a window of radius $H$ is at least $4\pi cH$;
using its explicit reciprocal derivative, integration by parts on each tail
bounds the difference between $G_\epsilon(c)$ and its damped window uniformly
in $\epsilon>0$ by $2/(cH\pi)$.
On the finite interval the damped integrand converges continuously to the
undamped one. The principal half-power is continuous at
$\pi/(2\pi c\ii)=-\ii/(2c)$, so
\[
 |F_c(H)-(-\ii/(2c))^{1/2}|\le2/(cH\pi).
\]
The number $(1-\ii)/(2\sqrt c)$ has positive real part and square
$-\ii/(2c)$, hence is that principal square root. Equivalently it is
$e^{-\ii\pi/4}/\sqrt{2c}$, proving the exact Fresnel normalization and its
finite-window error.

For the stationary integral, write $y=r+z$ with $r=r_b$ and
$\Theta(z)=\Phi(T,b;(r+z)^2)$.
For $|z|\le r/2$, the actual logarithm expansion gives
\[
 \Theta(z)=\Theta(0)-2\pi h_bz^2+\frac{2T}{3r^3}z^3+E_4(z),
 \qquad |E_4(z)|\le4T|z|^4/r^4.
\]
This follows from the logarithmic power-series remainder, multiplied by $2T$,
and the exact quadratic-root equation canceling its linear term.
If the root amplitude is $a(r+z)$, its $C^2$ bound gives
$|a(r+z)-a(r)-za'(r)|\le MR^2z^2$.
The cubic correction multiplied by $a(r)$ is odd, as is $za'(r)$; both
integrate to zero against the even quadratic exponential on $[-h,h]$.
The remaining terms are: the quadratic amplitude remainder;
the product of linear amplitude with cubic phase;
the fourth-order phase remainder times the bounded amplitude; and the
second-order exponential remainder of the cubic term.
Bounding these four terms before integration gives respectively
\[
 2MR^2h^3,\quad4MRT h^5/r^3,\quad8MTh^5/r^4,
                    \quad8MT^2h^7/r^6.
\]
This is the local symmetric estimate in Lemma~\ref{lem:stationaryerror}.
The source's tails use actual compact support, vanishing endpoint amplitudes,
and the nonstationary primitive estimates; the Fresnel-window bound then
provides the reciprocal $h$ term. Hence the fourth-root balance follows from
an exact odd-term cancellation, not from a generic stationary-phase theorem
with an unspecified remainder.
\src{Stronger/Atkinson/FresnelAbelLimit.lean; FresnelEvaluation.lean; StationaryOddRemainder.lean; AtkinsonSymmetricStationary.lean}{norm_atkinsonQuadraticWindow_sub_gaussianValue_le; fresnelGaussianValue_eq_phase; integral_atkinsonRootQuadratic_odd; IntervalC2Bound.atkinsonLocalSymmetric}

\subsection{Finite cancellation lemmas used twice in the development}
The Weyl producer and the Atkinson prefix packet both use the following finite
cancellation mechanism. It is proved locally in \lean{GuthMaynard}, not assumed
as an analytic input.
Let $\theta_n=f(n+1)-f(n)$ be monotone with
$\delta\le\theta_n\le2\pi-\delta$, $\delta>0$, $0\le n\le N$.
The phase resolvent
\[
 q(\theta)=\frac1{1-e^{\ii\theta}}
           =\frac12+\frac{\ii}2\cot(\theta/2)
\]
has constant real part. The identity
$e^{\ii f(n)}=q(\theta_n)(e^{\ii f(n)}-e^{\ii f(n+1)})$
telescopes by finite summation by parts. Its resulting norm bound is the two
endpoint resolvent norms plus the variation of $q(\theta_n)$.
Since $\cot$ is decreasing on $(0,\pi)$, that variation is
$(\cot(\theta_0/2)-\cot(\theta_N/2))/2$. Each endpoint satisfies
$|q(\theta)|\le\pi/(2\delta)$; its cotangent absolute value divided by two
is at most $|q(\theta)|$. Thus the phase sum is at most $2\pi/\delta$.
Subtracting $2\pi kn$ leaves every phase unchanged and reduces any fixed
period to this case. Replacing $f(n)$ by $-f(n)+2\pi n$ gives the decreasing
increment version by complex conjugation.

If instead the successive second differences satisfy
$0<\lambda\le\theta_{n+1}-\theta_n\le\Lambda$ for $n<N$, partition by
$k_n=\lfloor\theta_n/(2\pi)\rfloor$.
The index is monotone. The number of periods from $k_0$ to $k_N$ is at most
$N\Lambda/(2\pi)+2$, since $\theta_N-\theta_0\le N\Lambda$.
In each period the good increments lie in
$[2\pi k+\delta,2\pi(k+1)-\delta]$ and form a consecutive index set;
their phase sum is at most $2\pi/\delta$ by the resolvent estimate.
The two remaining endpoint bands have at most $\delta/\lambda+1$ indices
each: if $i$ and $j$ are in a band of width $\delta$, then
$(j-i)\lambda\le\theta_j-\theta_i\le\delta$.
Absolute estimation on those bands gives
\[
 \left|\sum_{n=0}^N e^{\ii f(n)}\right|
 \le\left(\frac{N\Lambda}{2\pi}+2\right)
                \left(\frac{2\pi}{\delta}+2(\delta/\lambda+1)\right).
\]
This estimate remains valid if the good band is empty or the two excluded
bands overlap, since the partition inequalities are upper bounds.
Taking $\delta=\sqrt\lambda$ proves the B-process used in
Section~\ref{sec:packets} with its exact constants.
\src{GuthMaynard/VanDerCorput.lean; GuthMaynard/SecondDerivative.lean}{kusminLandau_one_period; kusminLandau_one_period_decreasing; filter_card_bound; period_count_bound; vanDerCorput_second_derivative; vanDerCorput_B_process}

For completeness, the finite Weyl A-process is also an exact sum argument.
Extend a coefficient sequence by zero off $[0,N)$ and average its $H$ shifts
on $[-H,N)$. Each original term appears $H$ times.
Cauchy--Schwarz on this padded interval gives
\[
 H^2\left|\sum_{n<N}a_n\right|^2
 \le(N+H)\left(HN+H^2C\right)
\]
when $|a_n|\le1$ and every off-diagonal shifted correlation has norm at most
$C$. Expansion of the square supplies $HN$ from the $H$ diagonal shifts;
the other correlations have real part at most their norm, so their total is
at most $H^2C$.
For the logarithmic phase $f(x)=-t\log x$ on a block starting at $A$, the
shift difference has exact curvature
$t/u^2-t/v^2=tr(u+v)/(u^2v^2)$, where $v-u=r>0$ and
$A\le u\le v\le2A$. Hence its unit second differences lie between
$tr/(8A^3)$ and $4tr/A^3$. The preceding B-process and the elementary
product estimate, when $tr/(8A^3)\le1$, give the correlation bound
\[
 100\left(N\sqrt{tr/(8A^3)}+1/\sqrt{tr/(8A^3)}\right).
\]
Uniformize over $1\le r\le H$ by replacing the first square root with its
value at $H$ and the second denominator with its value at one.
For $t=Y^3$, $1\le Y\le A$, $A^2\le Y^3$, choose
$H=\lfloor A/Y\rfloor_+$. Then $1\le H\le A$,
$HY\le A\le2HY$. The uniform correlation is at most $400Y$.
With outer block length $A$, the A-process gives a square bound at most
$900AY$, hence phase-prefix norm at most $30\sqrt{AY}$. For shorter
prefixes, if their length is below $H$, use the triangle bound; otherwise
apply the same shift argument and enlarge its length factors to $A$.
Thus the prefix bound is genuinely uniform over all prefix lengths.

Finally for any nonnegative antitone weight $u_n$, summation by parts gives
\[
 \left|\sum_{n<N}u_nv_n\right|
 \le B\left(u_{N-1}+\sum_{n<N-1}(u_n-u_{n+1})\right)=Bu_0
\]
if every phase prefix has norm at most $B$. Take
$u_n=(A+n)^{-1/2}$ and $B=30\sqrt{AY}$ to get
$30\sqrt Y=30t^{1/6}$, the physical block estimate in
Proposition~\ref{prop:weyl}.
\src{GuthMaynard/Weyl.lean; GuthMaynard/WeylExplicit.lean; GuthMaynard/WeylZeta.lean; WeylPort/Block.lean}{interval_weyl_differencing_complex; logarithmicDifferenceCurvature_bounds; logarithmic_weyl_exponent_pair_prefix; norm_weighted_sum_le_of_antitone; norm_criticalLine_cpow_block_le}

\subsection{Summable nonstationary terms in the physical Bessel source}
Here are the concrete series majorants behind the removals in
Section~\ref{sec:atkinson}.
On the support $x\in[T/16,T]$, the actual $K_0$ integral has exponential
large-argument domination. For every fixed integer $k$, its pointwise
bound is at most $C_k/(T^kn^k)$ at $4\pi\sqrt{nx}$.
The source first uses $\cosh u\ge1+u^2/2$ to prove
$|K_0(z)|\le2/\sqrt z$ for $z>0$.
Splitting the exponential as two equal factors gives
$K_0(z)\le e^{-z/2}K_0(z/2)$, hence $|K_0(z)|\le2e^{-z/2}$ for $z\ge2$.
Since $e^{-z/2}\le(2k)!/(z/2)^{2k}$, and
$(4\pi\sqrt{nx})^2\ge Tn$ on the support, the power constant can be chosen
$C_k=2\cdot2^{2k}(2k)!$. The compact source test has integral norm at most
$CGT$. Therefore its $n$th full plus term is at most
\[
 \frac{C_kGT}{T^k}d(n)n^{-2}\qquad(k\ge2).
\]
The zero term vanishes. Since $\sum d(n)n^{-2}<\infty$, summing proves
$|\text{plus branch}|\le C_kGT/T^k$. Choose $k>A+2$ and then $T\ge C_k$
to obtain the power saving $GT^{-A}$.

The correction carriers in \eqref{eq:Ytwo} have $\alpha=3/4$.
The uniform root-kernel primitive bound is at most four, obtained by splitting
at $r_b-1,r_b+1$ and applying the one-sided nonstationary derivative bound
outside the middle interval of length two. Integration by parts with the
actual $C^1$ amplitude gives
$|\mathcal A_\alpha(T,G,L;b)|\le C_\alpha GT^{-\alpha}$ for every real $b$.
For $b\ge8\sqrt T$, the support is away from both signed saddles and the
stronger bound is $C_\alpha GT^{-\alpha}/b$ for each sign.
Combining these two bounds at $b=\sqrt n$, using
$T^{-3/4}\sqrt T\le1$, gives a correction pair bounded by $CG/\sqrt n$.
Its explicit Bessel factor $(4\pi)^{-3/2}n^{-3/4}/8$ and divisor factor
then give $CGd(n)n^{-5/4}$ per term. The convergent divisor series at
$5/4$ bounds the complete correction by $O(G)$.

The genuine remainder in \eqref{eq:Ytwo} is bounded on the same support by a
constant times $T^{-5/4}n^{-5/4}$. Multiply by the compact-test integral norm
$CGT$ and sum against $d(n)$; the resulting bound is
$CGT^{-1/4}\sum d(n)n^{-5/4}\le C'G$. These are full-series bounds,
so neither correction nor remainder can hide an unresolved infinite tail.

The logarithmic main term is written as a smooth weighted reflection integral
with kernel $x^{-1}e^{\ii(T\log x-2\pi x)}$ on $[T/16,T]$.
Its $C^1$ weight bound is $CG\sqrt T\log T$. The all-frequency estimate in
Lemma~\ref{lem:allfreq}, followed by integration by parts with this weight,
cancels the $\sqrt T$ scale and bounds the main term by $C'G\log T$.
This is why that term is absorbed into the local-mean error rather than set to
zero. The fixed support, coefficient norms, derivative integrals and primitive
estimates all enter before physical scales are specialized to $L=\log T$.
\src{Stronger/Atkinson/AtkinsonK0Integral.lean; AtkinsonK0Series.lean; AtkinsonCorrectionBounds.lean; NeumannRemainderSeries.lean; AtkinsonMainIntegral.lean}{exists_norm_integral_zetaAtkinsonK0_le; exists_norm_zetaAtkinsonBesselPlus_le; exists_norm_atkinsonCorrectionTerm_le; exists_norm_zetaAtkinsonVoronoiMain_le}

\section{External dependencies and the companion source ledger}
The pinned project uses Lean \texttt{v4.35.0-rc2}. Its mathlib revision is
specified in Section~\ref{sec:contract}. The other project dependencies
(plausible, LeanSearchClient, importGraph, proofwidgets, aesop, Qq,
batteries and Cli) are preserved by the recorded lake manifest and registry
record. They supply tooling and proof infrastructure; they are not unstated
analytic hypotheses.

The companion \texttt{source-ledger.md} resolves every displayed Lean
correspondence to the actual file and declaration starting line, gives a
SHA-256 for every local module in the textual import closure, and lists the
external import roots. The complete pinned source is archived beside it.
The ledger is a navigation and integrity record, not a substitute for the
proofs above. Its inventory is generated by text recognition and import
traversal, not by elaborating a Lean logical-dependency graph; generated
and unusually formatted declarations may not be counted. It does not claim
that each declaration in every imported library has been expanded into prose.
General imported foundations are used through their stated mathematical
interfaces, as is normal for an academic proof.



\end{document}
